\documentclass[8pt]{beamer}
\usetheme{Madrid}
\usecolortheme{default}
\usepackage{amsmath, amssymb, graphicx, array, multirow, booktabs, xcolor, tikz}
\usetikzlibrary{shapes,arrows,positioning,fit,backgrounds}

% Compact font settings
\setbeamerfont{title}{size=\Large}
\setbeamerfont{subtitle}{size=\small}
\setbeamerfont{author}{size=\scriptsize}
\setbeamerfont{date}{size=\scriptsize}
\setbeamerfont{frametitle}{size=\large}
\setbeamerfont{framesubtitle}{size=\normalsize}
\setbeamerfont{enumerate item}{size=\footnotesize}
\setbeamerfont{itemize item}{size=\footnotesize}
\setbeamerfont{block title}{size=\footnotesize}

% Compact spacing
\setlength{\itemsep}{0.08ex}
\setlength{\parskip}{0.08ex}
\setlength{\leftmargini}{0.3cm}
\setbeamersize{text margin left=3mm, text margin right=3mm}
\addtobeamertemplate{frame begin}{\vspace{-0.4cm}}{}

% Colors
\definecolor{myblue}{RGB}{0,0,255}
\definecolor{myred}{RGB}{255,0,0}
\definecolor{mygreen}{RGB}{0,128,0}
\definecolor{myorange}{RGB}{255,165,0}
\definecolor{mypurple}{RGB}{128,0,128}
\definecolor{mygray}{RGB}{128,128,128}
\definecolor{mygold}{RGB}{255,215,0}

\title{Introduction to Tools \& Techniques}
\subtitle{100 Questions: ML Fundamentals, Data Types, Pre-processing, PCA}
\author{GARP RAI Exam Preparation}
\date{}

\begin{document}
	
	\begin{frame}
		\titlepage
	\end{frame}
	
	% ============================================================
	% SECTION 1: MACHINE LEARNING FUNDAMENTALS - QUESTIONS 1-10
	% ============================================================
	
	\begin{frame}{Q1: Machine Learning Definition - Tutorial}
		\fontsize{6.5}{7.5}\selectfont
		\textbf{QUESTION: What is Machine Learning (ML)?}
		
		\vspace{0.1cm}
		\textbf{OPTIONS:}
		\begin{enumerate}
			\item A method for storing large datasets
			\item An umbrella term for techniques that train models to recognize data patterns for prediction and classification
			\item A software development framework
			\item A database management system
		\end{enumerate}
		
		\vspace{0.1cm}
		\textcolor{myblue}{\textbf{ANSWER: B}}
		
		\vspace{0.1cm}
		\textbf{DETAILED EXPLANATION:}
		\begin{itemize}
			\item ML is a subset of Artificial Intelligence (AI)
			\item Models are trained to recognize patterns in data
			\item Applications include prediction, classification, and risk management
			\item Techniques have evolved with advances in computing power and data availability
		\end{itemize}
		
		\vspace{0.1cm}
		\textcolor{myred}{\textbf{WHY OTHER OPTIONS ARE WRONG:}}
		\begin{itemize}
			\item \textcolor{myred}{A:} ML is about learning patterns, not just storage
			\item \textcolor{myred}{C:} ML is a set of techniques, not a software framework
			\item \textcolor{myred}{D:} ML is not a database system
		\end{itemize}
		
		\textcolor{myred}{\textbf{EXAM TRAP:}} ML = \textcolor{myblue}{training models to recognize patterns}!
	\end{frame}
	
	\begin{frame}{Q2: ML Applications - Tutorial}
		\fontsize{6.5}{7.5}\selectfont
		\textbf{QUESTION: Which is NOT an application of Machine Learning?}
		
		\vspace{0.1cm}
		\textbf{OPTIONS:}
		\begin{enumerate}
			\item Stock selection and trading
			\item Image recognition
			\item Manual data entry
			\item Fraud detection and credit scoring
		\end{enumerate}
		
		\vspace{0.1cm}
		\textcolor{myblue}{\textbf{ANSWER: C}}
		
		\vspace{0.1cm}
		\textbf{DETAILED EXPLANATION:}
		\begin{itemize}
			\item ML applications include stock selection, image recognition, game playing, autonomous vehicles
			\item Risk management: credit scoring, fraud detection
			\item Manual data entry is a human task, not an ML application
		\end{itemize}
		
		\vspace{0.1cm}
		\textcolor{myred}{\textbf{WHY OTHER OPTIONS ARE WRONG:}}
		\begin{itemize}
			\item \textcolor{myred}{A:} Stock selection is a common ML application
			\item \textcolor{myred}{B:} Image recognition is a major ML application
			\item \textcolor{myred}{D:} Fraud detection and credit scoring are common ML applications
		\end{itemize}
		
		\textcolor{myred}{\textbf{EXAM TRAP:}} ML applications include \textcolor{myblue}{prediction, classification, and pattern recognition}!
	\end{frame}
	
	\begin{frame}{Q3: ML Advantages - Tutorial}
		\fontsize{6.5}{7.5}\selectfont
		\textbf{QUESTION: What advantage does ML offer over traditional econometrics?}
		
		\vspace{0.1cm}
		\textbf{OPTIONS:}
		\begin{enumerate}
			\item It always produces more interpretable results
			\item It can handle big data, non-linearity, and dimensionality reduction
			\item It requires no data cleaning
			\item It never overfits
		\end{enumerate}
		
		\vspace{0.1cm}
		\textcolor{myblue}{\textbf{ANSWER: B}}
		
		\vspace{0.1cm}
		\textbf{DETAILED EXPLANATION:}
		\begin{itemize}
			\item \textbf{Handling big data:} ML scales better with large datasets
			\item \textbf{Handling non-linearity:} ML captures complex relationships
			\item \textbf{Reducing dimensionality:} ML techniques like PCA help with many features
			\item \textbf{Handling missing data:} ML has flexible imputation methods
		\end{itemize}
		
		\vspace{0.1cm}
		\textcolor{myred}{\textbf{WHY OTHER OPTIONS ARE WRONG:}}
		\begin{itemize}
			\item \textcolor{myred}{A:} ML models are often LESS interpretable
			\item \textcolor{myred}{C:} ML requires extensive data cleaning
			\item \textcolor{myred}{D:} ML models can overfit
		\end{itemize}
		
		\textcolor{myred}{\textbf{EXAM TRAP:}} ML advantages = \textcolor{myblue}{big data, non-linearity, dimensionality reduction}!
	\end{frame}
	
	\begin{frame}{Q4: ML vs Classical Statistics - Tutorial}
		\fontsize{6.5}{7.5}\selectfont
		\textbf{QUESTION: How does ML differ from classical statistics in model selection?}
		
		\vspace{0.1cm}
		\textbf{OPTIONS:}
		\begin{enumerate}
			\item ML relies on pre-specified theory; classical statistics lets data decide
			\item ML lets data decide; classical statistics relies on pre-specified theory
			\item Both rely equally on theory
			\item Neither relies on theory
		\end{enumerate}
		
		\vspace{0.1cm}
		\textcolor{myblue}{\textbf{ANSWER: B}}
		
		\vspace{0.1cm}
		\textbf{DETAILED EXPLANATION:}
		\begin{itemize}
			\item \textbf{Classical statistics:} Theory-driven, pre-specified model
			\item \textbf{ML:} Data-driven, uses techniques like regularization
			\item Classical: Hypothesis testing and inference
			\item ML: Focus on out-of-sample prediction
		\end{itemize}
		
		\vspace{0.1cm}
		\textcolor{myred}{\textbf{WHY OTHER OPTIONS ARE WRONG:}}
		\begin{itemize}
			\item \textcolor{myred}{A:} This reverses the roles
			\item \textcolor{myred}{C:} They differ fundamentally
			\item \textcolor{myred}{D:} Both rely on data, but differently
		\end{itemize}
		
		\textcolor{myred}{\textbf{EXAM TRAP:}} ML = \textcolor{myblue}{data-driven}; Classical = \textcolor{myblue}{theory-driven}!
	\end{frame}
	
	\begin{frame}{Q5: ML Focus - Tutorial}
		\fontsize{6.5}{7.5}\selectfont
		\textbf{QUESTION: What is the primary focus of Machine Learning?}
		
		\vspace{0.1cm}
		\textbf{OPTIONS:}
		\begin{enumerate}
			\item Inference and hypothesis testing
			\item Out-of-sample prediction accuracy
			\item R-squared maximization
			\item Coefficient interpretation
		\end{enumerate}
		
		\vspace{0.1cm}
		\textcolor{myblue}{\textbf{ANSWER: B}}
		
		\vspace{0.1cm}
		\textbf{DETAILED EXPLANATION:}
		\begin{itemize}
			\item ML focuses on predictive performance
			\item Out-of-sample accuracy is the key metric
			\item Generalization to new data is paramount
			\item Less emphasis on inference (p-values, t-statistics)
		\end{itemize}
		
		\vspace{0.1cm}
		\textcolor{myred}{\textbf{WHY OTHER OPTIONS ARE WRONG:}}
		\begin{itemize}
			\item \textcolor{myred}{A:} Inference is the focus of classical statistics
			\item \textcolor{myred}{C:} R-squared is a classical statistic
			\item \textcolor{myred}{D:} Coefficient interpretation is secondary in ML
		\end{itemize}
		
		\textcolor{myred}{\textbf{EXAM TRAP:}} ML focuses on \textcolor{myblue}{out-of-sample prediction}!
	\end{frame}
	
	\begin{frame}{Q6: ML Terminology - Tutorial}
		\fontsize{6.5}{7.5}\selectfont
		\textbf{QUESTION: What is the ML term for "dependent variable"?}
		
		\vspace{0.1cm}
		\textbf{OPTIONS:}
		\begin{enumerate}
			\item Feature
			\item Target, output, or label
			\item Instance
			\item Attribute
		\end{enumerate}
		
		\vspace{0.1cm}
		\textcolor{myblue}{\textbf{ANSWER: B}}
		
		\vspace{0.1cm}
		\textbf{DETAILED EXPLANATION:}
		\begin{itemize}
			\item \textbf{Statistics:} Dependent variable, regressand
			\item \textbf{ML:} Target, output, label, response variable
			\item In classification: label indicates category
			\item In prediction: target is what we predict
		\end{itemize}
		
		\vspace{0.1cm}
		\textcolor{myred}{\textbf{WHY OTHER OPTIONS ARE WRONG:}}
		\begin{itemize}
			\item \textcolor{myred}{A:} Feature is the independent variable
			\item \textcolor{myred}{C:} Instance is a data point/observation
			\item \textcolor{myred}{D:} Attribute is another term for feature
		\end{itemize}
		
		\textcolor{myred}{\textbf{EXAM TRAP:}} Dependent variable in ML = \textcolor{myblue}{target, output, or label}!
	\end{frame}
	
	\begin{frame}{Q7: ML Terminology (Independent Variable) - Tutorial}
		\fontsize{6.5}{7.5}\selectfont
		\textbf{QUESTION: What is the ML term for "independent variable"?}
		
		\vspace{0.1cm}
		\textbf{OPTIONS:}
		\begin{enumerate}
			\item Output
			\item Label
			\item Feature, input, or signal
			\item Instance
		\end{enumerate}
		
		\vspace{0.1cm}
		\textcolor{myblue}{\textbf{ANSWER: C}}
		
		\vspace{0.1cm}
		\textbf{DETAILED EXPLANATION:}
		\begin{itemize}
			\item \textbf{Statistics:} Independent variable, predictor, regressor
			\item \textbf{ML:} Feature, input, signal, attribute
			\item Used to predict the target/output
			\item Can be numerical or categorical
		\end{itemize}
		
		\vspace{0.1cm}
		\textcolor{myred}{\textbf{WHY OTHER OPTIONS ARE WRONG:}}
		\begin{itemize}
			\item \textcolor{myred}{A:} Output is the dependent variable
			\item \textcolor{myred}{B:} Label is the dependent variable in classification
			\item \textcolor{myred}{D:} Instance is a data point
		\end{itemize}
		
		\textcolor{myred}{\textbf{EXAM TRAP:}} Independent variable in ML = \textcolor{myblue}{feature, input, or signal}!
	\end{frame}
	
	\begin{frame}{Q8: ML Terminology (Estimation) - Tutorial}
		\fontsize{6.5}{7.5}\selectfont
		\textbf{QUESTION: What is the ML term for "estimation"?}
		
		\vspace{0.1cm}
		\textbf{OPTIONS:}
		\begin{enumerate}
			\item Hypothesis testing
			\item Training or learning
			\item Model selection
			\item Validation
		\end{enumerate}
		
		\vspace{0.1cm}
		\textcolor{myblue}{\textbf{ANSWER: B}}
		
		\vspace{0.1cm}
		\textbf{DETAILED EXPLANATION:}
		\begin{itemize}
			\item \textbf{Statistics:} Estimation
			\item \textbf{ML:} Training or learning
			\item The estimator is called a learner or algorithm
			\item In classification: called a classifier
		\end{itemize}
		
		\vspace{0.1cm}
		\textcolor{myred}{\textbf{WHY OTHER OPTIONS ARE WRONG:}}
		\begin{itemize}
			\item \textcolor{myred}{A:} Hypothesis testing is classical statistics
			\item \textcolor{myred}{C:} Model selection is separate from estimation
			\item \textcolor{myred}{D:} Validation is for model evaluation
		\end{itemize}
		
		\textcolor{myred}{\textbf{EXAM TRAP:}} Estimation in ML = \textcolor{myblue}{training or learning}!
	\end{frame}
	
	\begin{frame}{Q9: Data Point in ML - Tutorial}
		\fontsize{6.5}{7.5}\selectfont
		\textbf{QUESTION: What is the ML term for a "data point"?}
		
		\vspace{0.1cm}
		\textbf{OPTIONS:}
		\begin{enumerate}
			\item Feature
			\item Label
			\item Instance or example
			\item Target
		\end{enumerate}
		
		\vspace{0.1cm}
		\textcolor{myblue}{\textbf{ANSWER: C}}
		
		\vspace{0.1cm}
		\textbf{DETAILED EXPLANATION:}
		\begin{itemize}
			\item \textbf{Statistics:} Observation, data point
			\item \textbf{ML:} Instance, example
			\item Each instance has features and possibly a label
			\item Represents one row in the dataset
		\end{itemize}
		
		\vspace{0.1cm}
		\textcolor{myred}{\textbf{WHY OTHER OPTIONS ARE WRONG:}}
		\begin{itemize}
			\item \textcolor{myred}{A:} Feature is a column/variable
			\item \textcolor{myred}{B:} Label is the target/outcome
			\item \textcolor{myred}{D:} Target is the dependent variable
		\end{itemize}
		
		\textcolor{myred}{\textbf{EXAM TRAP:}} Data point in ML = \textcolor{myblue}{instance or example}!
	\end{frame}
	
	\begin{frame}{Q10: Terminology Summary - Tutorial}
		\fontsize{6.5}{7.5}\selectfont
		\textbf{QUESTION: Match the ML term with its statistics equivalent: "Feature"}
		
		\vspace{0.1cm}
		\textbf{OPTIONS:}
		\begin{enumerate}
			\item Dependent variable
			\item Independent variable or predictor
			\item Data point
			\item Estimation
		\end{enumerate}
		
		\vspace{0.1cm}
		\textcolor{myblue}{\textbf{ANSWER: B}}
		
		\vspace{0.1cm}
		\textbf{DETAILED EXPLANATION:}
		\begin{itemize}
			\item \textbf{Statistics:} Independent variable, predictor, regressor
			\item \textbf{ML:} Feature, input, signal, attribute
			\item Used to explain/predict the target
			\item Can be continuous or categorical
		\end{itemize}
		
		\vspace{0.1cm}
		\textcolor{myred}{\textbf{WHY OTHER OPTIONS ARE WRONG:}}
		\begin{itemize}
			\item \textcolor{myred}{A:} This is Target/Label in ML
			\item \textcolor{myred}{C:} This is Instance/Example in ML
			\item \textcolor{myred}{D:} This is Training/Learning in ML
		\end{itemize}
		
		\textcolor{myred}{\textbf{EXAM TRAP:}} Feature = \textcolor{myblue}{independent variable} in statistics!
	\end{frame}
	
	% ============================================================
	% SECTION 2: TYPES OF MACHINE LEARNING - QUESTIONS 11-25
	% ============================================================
	
	\begin{frame}{Q11: Unsupervised Learning Definition - Tutorial}
		\fontsize{6.5}{7.5}\selectfont
		\textbf{QUESTION: What is unsupervised learning?}
		
		\vspace{0.1cm}
		\textbf{OPTIONS:}
		\begin{enumerate}
			\item Learning with labeled data to predict outcomes
			\item Learning with unlabeled data to recognize patterns
			\item Learning through trial and error with rewards
			\item Learning with partially labeled data
		\end{enumerate}
		
		\vspace{0.1cm}
		\textcolor{myblue}{\textbf{ANSWER: B}}
		
		\vspace{0.1cm}
		\textbf{DETAILED EXPLANATION:}
		\begin{itemize}
			\item No corresponding output/label for each observation
			\item Goal: recognize patterns, understand structure
			\item Applications: clustering, anomaly detection
			\item Examples: PCA, K-means clustering
		\end{itemize}
		
		\vspace{0.1cm}
		\textcolor{myred}{\textbf{WHY OTHER OPTIONS ARE WRONG:}}
		\begin{itemize}
			\item \textcolor{myred}{A:} This describes supervised learning
			\item \textcolor{myred}{C:} This describes reinforcement learning
			\item \textcolor{myred}{D:} This describes semi-supervised learning
		\end{itemize}
		
		\textcolor{myred}{\textbf{EXAM TRAP:}} Unsupervised = \textcolor{myblue}{unlabeled data, pattern recognition}!
	\end{frame}
	
	\begin{frame}{Q12: Unsupervised Learning Applications - Tutorial}
		\fontsize{6.5}{7.5}\selectfont
		\textbf{QUESTION: Which is an application of unsupervised learning?}
		
		\vspace{0.1cm}
		\textbf{OPTIONS:}
		\begin{enumerate}
			\item Predicting house prices
			\item Clustering stocks with similar characteristics
			\item Classifying loan defaults
			\item Forecasting stock returns
		\end{enumerate}
		
		\vspace{0.1cm}
		\textcolor{myblue}{\textbf{ANSWER: B}}
		
		\vspace{0.1cm}
		\textbf{DETAILED EXPLANATION:}
		\begin{itemize}
			\item Clustering groups similar items without labels
			\item Other applications: anomaly detection, dimensionality reduction (PCA)
			\item Used to understand data structure
			\item Helps identify patterns before supervised learning
		\end{itemize}
		
		\vspace{0.1cm}
		\textcolor{myred}{\textbf{WHY OTHER OPTIONS ARE WRONG:}}
		\begin{itemize}
			\item \textcolor{myred}{A:} Predicting house prices is supervised (regression)
			\item \textcolor{myred}{C:} Classifying defaults is supervised (classification)
			\item \textcolor{myred}{D:} Forecasting returns is supervised (time series prediction)
		\end{itemize}
		
		\textcolor{myred}{\textbf{EXAM TRAP:}} Unsupervised = \textcolor{myblue}{clustering and pattern discovery}!
	\end{frame}
	
	\begin{frame}{Q13: Supervised Learning Definition - Tutorial}
		\fontsize{6.5}{7.5}\selectfont
		\textbf{QUESTION: What is supervised learning?}
		
		\vspace{0.1cm}
		\textbf{OPTIONS:}
		\begin{enumerate}
			\item Learning with unlabeled data
			\item Learning with labeled data to predict outcomes
			\item Learning through rewards and punishments
			\item Learning with partially labeled data
		\end{enumerate}
		
		\vspace{0.1cm}
		\textcolor{myblue}{\textbf{ANSWER: B}}
		
		\vspace{0.1cm}
		\textbf{DETAILED EXPLANATION:}
		\begin{itemize}
			\item Each observation has features AND a label/output
			\item Goal: predict target for new, unseen data
			\item Two types: prediction (continuous) and classification (categorical)
			\item Examples: linear regression, logistic regression
		\end{itemize}
		
		\vspace{0.1cm}
		\textcolor{myred}{\textbf{WHY OTHER OPTIONS ARE WRONG:}}
		\begin{itemize}
			\item \textcolor{myred}{A:} This describes unsupervised learning
			\item \textcolor{myred}{C:} This describes reinforcement learning
			\item \textcolor{myred}{D:} This describes semi-supervised learning
		\end{itemize}
		
		\textcolor{myred}{\textbf{EXAM TRAP:}} Supervised = \textcolor{myblue}{labeled data, prediction/classification}!
	\end{frame}
	
	\begin{frame}{Q14: Prediction vs Classification - Tutorial}
		\fontsize{6.5}{7.5}\selectfont
		\textbf{QUESTION: What is the difference between prediction and classification?}
		
		\vspace{0.1cm}
		\textbf{OPTIONS:}
		\begin{enumerate}
			\item Prediction uses categorical variables; classification uses continuous
			\item Prediction predicts numerical values; classification assigns categories
			\item Prediction is unsupervised; classification is supervised
			\item There is no difference
		\end{enumerate}
		
		\vspace{0.1cm}
		\textcolor{myblue}{\textbf{ANSWER: B}}
		
		\vspace{0.1cm}
		\textbf{DETAILED EXPLANATION:}
		\begin{itemize}
			\item \textbf{Prediction:} Numerical output (e.g., house price)
			\item \textbf{Classification:} Categorical output (e.g., default/no default)
			\item Both are forms of supervised learning
			\item Different evaluation metrics apply
		\end{itemize}
		
		\vspace{0.1cm}
		\textcolor{myred}{\textbf{WHY OTHER OPTIONS ARE WRONG:}}
		\begin{itemize}
			\item \textcolor{myred}{A:} This is reversed
			\item \textcolor{myred}{C:} Both are supervised
			\item \textcolor{myred}{D:} They are different types of supervised learning
		\end{itemize}
		
		\textcolor{myred}{\textbf{EXAM TRAP:}} Prediction = \textcolor{myblue}{numerical}; Classification = \textcolor{myblue}{categorical}!
	\end{frame}
	
	\begin{frame}{Q15: Classification Example - Tutorial}
		\fontsize{6.5}{7.5}\selectfont
		\textbf{QUESTION: Which is an example of a classification problem?}
		
		\vspace{0.1cm}
		\textbf{OPTIONS:}
		\begin{enumerate}
			\item Forecasting GDP growth
			\item Predicting stock returns
			\item Determining if a loan will default
			\item Estimating house prices
		\end{enumerate}
		
		\vspace{0.1cm}
		\textcolor{myblue}{\textbf{ANSWER: C}}
		
		\vspace{0.1cm}
		\textbf{DETAILED EXPLANATION:}
		\begin{itemize}
			\item Classification: assigning to categories
			\item Loan default: Yes/No (binary classification)
			\item Target variable is categorical
			\item Other examples: fraud detection, credit rating
		\end{itemize}
		
		\vspace{0.1cm}
		\textcolor{myred}{\textbf{WHY OTHER OPTIONS ARE WRONG:}}
		\begin{itemize}
			\item \textcolor{myred}{A:} GDP growth is numerical (prediction/regression)
			\item \textcolor{myred}{B:} Stock returns are numerical (prediction)
			\item \textcolor{myred}{D:} House prices are numerical (prediction)
		\end{itemize}
		
		\textcolor{myred}{\textbf{EXAM TRAP:}} Classification = \textcolor{myblue}{categorical output} like default/no default!
	\end{frame}
	
	\begin{frame}{Q16: Prediction Example - Tutorial}
		\fontsize{6.5}{7.5}\selectfont
		\textbf{QUESTION: Which is an example of a prediction problem?}
		
		\vspace{0.1cm}
		\textbf{OPTIONS:}
		\begin{enumerate}
			\item Email spam detection
			\item Loan default classification
			\item Predicting house sale prices
			\item Image recognition
		\end{enumerate}
		
		\vspace{0.1cm}
		\textcolor{myblue}{\textbf{ANSWER: C}}
		
		\vspace{0.1cm}
		\textbf{DETAILED EXPLANATION:}
		\begin{itemize}
			\item Prediction: continuous numerical output
			\item House price: continuous numerical value
			\item Can be cross-sectional or time-series
			\item Examples: stock returns, GDP, revenue
		\end{itemize}
		
		\vspace{0.1cm}
		\textcolor{myred}{\textbf{WHY OTHER OPTIONS ARE WRONG:}}
		\begin{itemize}
			\item \textcolor{myred}{A:} Email spam is classification (spam/not spam)
			\item \textcolor{myred}{B:} Loan default is classification (default/no default)
			\item \textcolor{myred}{D:} Image recognition is classification (object categories)
		\end{itemize}
		
		\textcolor{myred}{\textbf{EXAM TRAP:}} Prediction = \textcolor{myblue}{continuous numerical output}!
	\end{frame}
	
	\begin{frame}{Q17: Semi-Supervised Learning - Tutorial}
		\fontsize{6.5}{7.5}\selectfont
		\textbf{QUESTION: What is semi-supervised learning?}
		
		\vspace{0.1cm}
		\textbf{OPTIONS:}
		\begin{enumerate}
			\item All data is labeled
			\item No data is labeled
			\item Only part of the data is labeled
			\item Learning through rewards and punishments
		\end{enumerate}
		
		\vspace{0.1cm}
		\textcolor{myblue}{\textbf{ANSWER: C}}
		
		\vspace{0.1cm}
		\textbf{DETAILED EXPLANATION:}
		\begin{itemize}
			\item Mix of labeled and unlabeled data
			\item Uses labeled data for supervised learning
			\item Uses unlabeled data to learn patterns in features
			\item Can improve performance with limited labeled data
		\end{itemize}
		
		\vspace{0.1cm}
		\textcolor{myred}{\textbf{WHY OTHER OPTIONS ARE WRONG:}}
		\begin{itemize}
			\item \textcolor{myred}{A:} This is supervised learning
			\item \textcolor{myred}{B:} This is unsupervised learning
			\item \textcolor{myred}{D:} This is reinforcement learning
		\end{itemize}
		
		\textcolor{myred}{\textbf{EXAM TRAP:}} Semi-supervised = \textcolor{myblue}{partially labeled data}!
	\end{frame}
	
	\begin{frame}{Q18: Reinforcement Learning Definition - Tutorial}
		\fontsize{6.5}{7.5}\selectfont
		\textbf{QUESTION: What is reinforcement learning?}
		
		\vspace{0.1cm}
		\textbf{OPTIONS:}
		\begin{enumerate}
			\item Learning from labeled examples
			\item Learning from unlabeled data
			\item Learning through trial and error with rewards
			\item Learning with partially labeled data
		\end{enumerate}
		
		\vspace{0.1cm}
		\textcolor{myblue}{\textbf{ANSWER: C}}
		
		\vspace{0.1cm}
		\textbf{DETAILED EXPLANATION:}
		\begin{itemize}
			\item No explicit labels
			\item Feedback through rewards/penalties
			\item Trial-and-error approach
			\item Output = recommended action given circumstances
		\end{itemize}
		
		\vspace{0.1cm}
		\textcolor{myred}{\textbf{WHY OTHER OPTIONS ARE WRONG:}}
		\begin{itemize}
			\item \textcolor{myred}{A:} This is supervised learning
			\item \textcolor{myred}{B:} This is unsupervised learning
			\item \textcolor{myred}{D:} This is semi-supervised learning
		\end{itemize}
		
		\textcolor{myred}{\textbf{EXAM TRAP:}} Reinforcement = \textcolor{myblue}{trial and error with rewards}!
	\end{frame}
	
	\begin{frame}{Q19: Reinforcement Learning Applications - Tutorial}
		\fontsize{6.5}{7.5}\selectfont
		\textbf{QUESTION: Which is an application of reinforcement learning in finance?}
		
		\vspace{0.1cm}
		\textbf{OPTIONS:}
		\begin{enumerate}
			\item Customer segmentation
			\item Optimal portfolio management
			\item Fraud detection
			\item House price prediction
		\end{enumerate}
		
		\vspace{0.1cm}
		\textcolor{myblue}{\textbf{ANSWER: B}}
		
		\vspace{0.1cm}
		\textbf{DETAILED EXPLANATION:}
		\begin{itemize}
			\item Reinforcement learning: making sequential decisions
			\item Portfolio management: buy/sell decisions over time
			\item Also: optimal trading, hedging, execution
			\item Learns from rewards (returns, risk reduction)
		\end{itemize}
		
		\vspace{0.1cm}
		\textcolor{myred}{\textbf{WHY OTHER OPTIONS ARE WRONG:}}
		\begin{itemize}
			\item \textcolor{myred}{A:} Customer segmentation is unsupervised
			\item \textcolor{myred}{C:} Fraud detection is supervised classification
			\item \textcolor{myred}{D:} House price prediction is supervised regression
		\end{itemize}
		
		\textcolor{myred}{\textbf{EXAM TRAP:}} Reinforcement = \textcolor{myblue}{sequential decision making} like portfolio management!
	\end{frame}
	
	\begin{frame}{Q20: Parametric Methods - Tutorial}
		\fontsize{6.5}{7.5}\selectfont
		\textbf{QUESTION: What characterizes parametric machine learning methods?}
		
		\vspace{0.1cm}
		\textbf{OPTIONS:}
		\begin{enumerate}
			\item No assumptions about functional form
			\item Make assumptions about functional form between features and output
			\item Always outperform non-parametric methods
			\item Require no training data
		\end{enumerate}
		
		\vspace{0.1cm}
		\textcolor{myblue}{\textbf{ANSWER: B}}
		
		\vspace{0.1cm}
		\textbf{DETAILED EXPLANATION:}
		\begin{itemize}
			\item Assume a specific functional form (e.g., linear)
			\item Parameters are estimated from data
			\item Risk: wrong functional form = poor fit
			\item Examples: linear regression, logistic regression
		\end{itemize}
		
		\vspace{0.1cm}
		\textcolor{myred}{\textbf{WHY OTHER OPTIONS ARE WRONG:}}
		\begin{itemize}
			\item \textcolor{myred}{A:} This describes non-parametric methods
			\item \textcolor{myred}{C:} Not always true
			\item \textcolor{myred}{D:} Parametric methods require training data
		\end{itemize}
		
		\textcolor{myred}{\textbf{EXAM TRAP:}} Parametric = \textcolor{myblue}{assumes functional form}!
	\end{frame}
	
	\begin{frame}{Q21: Non-Parametric Methods - Tutorial}
		\fontsize{6.5}{7.5}\selectfont
		\textbf{QUESTION: What characterizes non-parametric machine learning methods?}
		
		\vspace{0.1cm}
		\textbf{OPTIONS:}
		\begin{enumerate}
			\item Make assumptions about functional form
			\item No explicit assumption about functional form
			\item Require fewer observations
			\item Are always less accurate
		\end{enumerate}
		
		\vspace{0.1cm}
		\textcolor{myblue}{\textbf{ANSWER: B}}
		
		\vspace{0.1cm}
		\textbf{DETAILED EXPLANATION:}
		\begin{itemize}
			\item Flexible - no fixed functional form
			\item Capture complex patterns
			\item Require many observations for accuracy
			\item Examples: decision trees, KNN, neural networks
		\end{itemize}
		
		\vspace{0.1cm}
		\textcolor{myred}{\textbf{WHY OTHER OPTIONS ARE WRONG:}}
		\begin{itemize}
			\item \textcolor{myred}{A:} This describes parametric methods
			\item \textcolor{myred}{C:} Non-parametric typically requires MORE observations
			\item \textcolor{myred}{D:} Can be very accurate with enough data
		\end{itemize}
		
		\textcolor{myred}{\textbf{EXAM TRAP:}} Non-parametric = \textcolor{myblue}{no fixed functional form}!
	\end{frame}
	
	\begin{frame}{Q22: Parametric vs Non-Parametric - Tutorial}
		\fontsize{6.5}{7.5}\selectfont
		\textbf{QUESTION: What is a disadvantage of parametric methods?}
		
		\vspace{0.1cm}
		\textbf{OPTIONS:}
		\begin{enumerate}
			\item They are too flexible
			\item Risk of wrong functional form leading to poor fit
			\item They require too much data
			\item They are always slow
		\end{enumerate}
		
		\vspace{0.1cm}
		\textcolor{myblue}{\textbf{ANSWER: B}}
		
		\vspace{0.1cm}
		\textbf{DETAILED EXPLANATION:}
		\begin{itemize}
			\item Parametric methods assume a specific form
			\item If assumption is wrong, model fails to fit data
			\item Example: linear model for quadratic relationship
			\item Non-parametric methods avoid this risk
		\end{itemize}
		
		\vspace{0.1cm}
		\textcolor{myred}{\textbf{WHY OTHER OPTIONS ARE WRONG:}}
		\begin{itemize}
			\item \textcolor{myred}{A:} Non-parametric are more flexible
			\item \textcolor{myred}{C:} Non-parametric require more data
			\item \textcolor{myred}{D:} Speed depends on implementation
		\end{itemize}
		
		\textcolor{myred}{\textbf{EXAM TRAP:}} Parametric risk = \textcolor{myblue}{wrong functional form}!
	\end{frame}
	
	\begin{frame}{Q23: Non-Parametric Requirement - Tutorial}
		\fontsize{6.5}{7.5}\selectfont
		\textbf{QUESTION: What is a requirement for non-parametric methods?}
		
		\vspace{0.1cm}
		\textbf{OPTIONS:}
		\begin{enumerate}
			\item Few observations
			\item Many observations for accurate estimation
			\item No data at all
			\item Pre-specified functional form
		\end{enumerate}
		
		\vspace{0.1cm}
		\textcolor{myblue}{\textbf{ANSWER: B}}
		
		\vspace{0.1cm}
		\textbf{DETAILED EXPLANATION:}
		\begin{itemize}
			\item Non-parametric methods are data-intensive
			\item Need sufficient data to capture complex patterns
			\item More data = better estimation
			\item Parametric methods can work with less data
		\end{itemize}
		
		\vspace{0.1cm}
		\textcolor{myred}{\textbf{WHY OTHER OPTIONS ARE WRONG:}}
		\begin{itemize}
			\item \textcolor{myred}{A:} Non-parametric needs MANY observations
			\item \textcolor{myred}{C:} All methods need data
			\item \textcolor{myred}{D:} This describes parametric methods
		\end{itemize}
		
		\textcolor{myred}{\textbf{EXAM TRAP:}} Non-parametric = \textcolor{myblue}{needs many observations}!
	\end{frame}
	
	\begin{frame}{Q24: ML Categories Summary - Tutorial}
		\fontsize{6.5}{7.5}\selectfont
		\textbf{QUESTION: What are the four main categories of machine learning?}
		
		\vspace{0.1cm}
		\textbf{OPTIONS:}
		\begin{enumerate}
			\item Supervised, Unsupervised, Semi-supervised, Reinforcement
			\item Training, Validation, Testing, Deployment
			\item Regression, Classification, Clustering, Dimensionality Reduction
			\item Parametric, Non-parametric, Linear, Non-linear
		\end{enumerate}
		
		\vspace{0.1cm}
		\textcolor{myblue}{\textbf{ANSWER: A}}
		
		\vspace{0.1cm}
		\textbf{DETAILED EXPLANATION:}
		\begin{itemize}
			\item \textbf{Supervised:} Labeled data, prediction/classification
			\item \textbf{Unsupervised:} Unlabeled data, pattern discovery
			\item \textbf{Semi-supervised:} Partially labeled data
			\item \textbf{Reinforcement:} Trial and error, rewards
		\end{itemize}
		
		\vspace{0.1cm}
		\textcolor{myred}{\textbf{WHY OTHER OPTIONS ARE WRONG:}}
		\begin{itemize}
			\item \textcolor{myred}{B:} These are data splitting stages
			\item \textcolor{myred}{C:} These are types of problems/tasks
			\item \textcolor{myred}{D:} These are modeling approaches
		\end{itemize}
		
		\textcolor{myred}{\textbf{EXAM TRAP:}} Four categories = \textcolor{myblue}{Supervised, Unsupervised, Semi-supervised, Reinforcement}!
	\end{frame}
	
	\begin{frame}{Q25: ML Data Requirements - Tutorial}
		\fontsize{6.5}{7.5}\selectfont
		\textbf{QUESTION: Which ML type requires labeled data?}
		
		\vspace{0.1cm}
		\textbf{OPTIONS:}
		\begin{enumerate}
			\item Unsupervised learning
			\item Supervised learning
			\item Reinforcement learning
			\item Both supervised and semi-supervised
		\end{enumerate}
		
		\vspace{0.1cm}
		\textcolor{myblue}{\textbf{ANSWER: D}}
		
		\vspace{0.1cm}
		\textbf{DETAILED EXPLANATION:}
		\begin{itemize}
			\item Supervised: fully labeled data
			\item Semi-supervised: partially labeled data
			\item Unsupervised: no labels
			\item Reinforcement: no labels, only rewards
		\end{itemize}
		
		\vspace{0.1cm}
		\textcolor{myred}{\textbf{WHY OTHER OPTIONS ARE WRONG:}}
		\begin{itemize}
			\item \textcolor{myred}{A:} Unsupervised has NO labels
			\item \textcolor{myred}{B:} Correct but incomplete
			\item \textcolor{myred}{C:} Reinforcement has no labels, only rewards
		\end{itemize}
		
		\textcolor{myred}{\textbf{EXAM TRAP:}} Supervised and Semi-supervised = \textcolor{myblue}{labeled data}!
	\end{frame}
	
	% ============================================================
	% SECTION 3: DATA TYPES - QUESTIONS 26-40
	% ============================================================
	
	\begin{frame}{Q26: Structured Data - Tutorial}
		\fontsize{6.5}{7.5}\selectfont
		\textbf{QUESTION: What is structured data?}
		
		\vspace{0.1cm}
		\textbf{OPTIONS:}
		\begin{enumerate}
			\item Data with no organization
			\item Data organized in rows and columns
			\item Data from images and videos
			\item Data in natural language
		\end{enumerate}
		
		\vspace{0.1cm}
		\textcolor{myblue}{\textbf{ANSWER: B}}
		
		\vspace{0.1cm}
		\textbf{DETAILED EXPLANATION:}
		\begin{itemize}
			\item Organized in rows and columns
			\item Each column = attribute/feature
			\item Each row = observation/instance
			\item Examples: census data, bank customer data
		\end{itemize}
		
		\vspace{0.1cm}
		\textcolor{myred}{\textbf{WHY OTHER OPTIONS ARE WRONG:}}
		\begin{itemize}
			\item \textcolor{myred}{A:} This is unstructured data
			\item \textcolor{myred}{C:} This is unstructured data (images/videos)
			\item \textcolor{myred}{D:} This is unstructured data (text)
		\end{itemize}
		
		\textcolor{myred}{\textbf{EXAM TRAP:}} Structured data = \textcolor{myblue}{rows and columns}!
	\end{frame}
	
	\begin{frame}{Q27: Unstructured Data - Tutorial}
		\fontsize{6.5}{7.5}\selectfont
		\textbf{QUESTION: Which is an example of unstructured data?}
		
		\vspace{0.1cm}
		\textbf{OPTIONS:}
		\begin{enumerate}
			\item Census data
			\item Bank customer database
			\item Social media text posts
			\item Excel spreadsheet with numeric data
		\end{enumerate}
		
		\vspace{0.1cm}
		\textcolor{myblue}{\textbf{ANSWER: C}}
		
		\vspace{0.1cm}
		\textbf{DETAILED EXPLANATION:}
		\begin{itemize}
			\item Unstructured: not arranged in rows/columns
			\item Examples: text, images, audio, video
			\item Harder to analyze
			\item Requires pre-processing (NLP for text)
		\end{itemize}
		
		\vspace{0.1cm}
		\textcolor{myred}{\textbf{WHY OTHER OPTIONS ARE WRONG:}}
		\begin{itemize}
			\item \textcolor{myred}{A:} Census data is structured
			\item \textcolor{myred}{B:} Bank database is structured
			\item \textcolor{myred}{D:} Excel with numeric data is structured
		\end{itemize}
		
		\textcolor{myred}{\textbf{EXAM TRAP:}} Unstructured = \textcolor{myblue}{text, images, audio, video}!
	\end{frame}
	
	\begin{frame}{Q28: Semi-Structured Data - Tutorial}
		\fontsize{6.5}{7.5}\selectfont
		\textbf{QUESTION: What is semi-structured data?}
		
		\vspace{0.1cm}
		\textbf{OPTIONS:}
		\begin{enumerate}
			\item Data with no structure at all
			\item Data with some structure but not fixed format
			\item Data in rows and columns
			\item Data that is always numeric
		\end{enumerate}
		
		\vspace{0.1cm}
		\textcolor{myblue}{\textbf{ANSWER: B}}
		
		\vspace{0.1cm}
		\textbf{DETAILED EXPLANATION:}
		\begin{itemize}
			\item Partly structured, partly unstructured
			\item Has some organization but not rigid
			\item Examples: emails, CSV files, HTML files
			\item Can contain both structured and unstructured elements
		\end{itemize}
		
		\vspace{0.1cm}
		\textcolor{myred}{\textbf{WHY OTHER OPTIONS ARE WRONG:}}
		\begin{itemize}
			\item \textcolor{myred}{A:} This is unstructured data
			\item \textcolor{myred}{C:} This is structured data
			\item \textcolor{myred}{D:} Can include text as well
		\end{itemize}
		
		\textcolor{myred}{\textbf{EXAM TRAP:}} Semi-structured = \textcolor{myblue}{some structure, not fixed format}!
	\end{frame}
	
	\begin{frame}{Q29: Numerical Data - Tutorial}
		\fontsize{6.5}{7.5}\selectfont
		\textbf{QUESTION: What is numerical data?}
		
		\vspace{0.1cm}
		\textbf{OPTIONS:}
		\begin{enumerate}
			\item Data with categories
			\item Data with numerical values that have natural ordering
			\item Data from images
			\item Data from text
		\end{enumerate}
		
		\vspace{0.1cm}
		\textcolor{myblue}{\textbf{ANSWER: B}}
		
		\vspace{0.1cm}
		\textbf{DETAILED EXPLANATION:}
		\begin{itemize}
			\item Continuous or discrete values
			\item Has natural ordering (\$50,000 > \$40,000)
			\item Can be used directly in ML models
			\item Examples: age, income, temperature
		\end{itemize}
		
		\vspace{0.1cm}
		\textcolor{myred}{\textbf{WHY OTHER OPTIONS ARE WRONG:}}
		\begin{itemize}
			\item \textcolor{myred}{A:} This is categorical data
			\item \textcolor{myred}{C:} This is unstructured data
			\item \textcolor{myred}{D:} This is unstructured data
		\end{itemize}
		
		\textcolor{myred}{\textbf{EXAM TRAP:}} Numerical data = \textcolor{myblue}{values with natural ordering}!
	\end{frame}
	
	\begin{frame}{Q30: Categorical Data - Tutorial}
		\fontsize{6.5}{7.5}\selectfont
		\textbf{QUESTION: What is categorical data?}
		
		\vspace{0.1cm}
		\textbf{OPTIONS:}
		\begin{enumerate}
			\item Data with continuous values
			\item Data that takes discrete values representing categories
			\item Data that is always numeric
			\item Data with no structure
		\end{enumerate}
		
		\vspace{0.1cm}
		\textcolor{myblue}{\textbf{ANSWER: B}}
		
		\vspace{0.1cm}
		\textbf{DETAILED EXPLANATION:}
		\begin{itemize}
			\item Takes discrete values
			\item Represents groups/categories
			\item Can be nominal (no order) or ordinal (ordered)
			\item Examples: marital status, gender, region
		\end{itemize}
		
		\vspace{0.1cm}
		\textcolor{myred}{\textbf{WHY OTHER OPTIONS ARE WRONG:}}
		\begin{itemize}
			\item \textcolor{myred}{A:} This is numerical data
			\item \textcolor{myred}{C:} Categorical can be non-numeric
			\item \textcolor{myred}{D:} This is unstructured data
		\end{itemize}
		
		\textcolor{myred}{\textbf{EXAM TRAP:}} Categorical data = \textcolor{myblue}{discrete categories}!
	\end{frame}
	
	\begin{frame}{Q31: Nominal Data - Tutorial}
		\fontsize{6.5}{7.5}\selectfont
		\textbf{QUESTION: What is nominal data?}
		
		\vspace{0.1cm}
		\textbf{OPTIONS:}
		\begin{enumerate}
			\item Categorical data with natural ordering
			\item Categorical data with no natural ordering
			\item Continuous numerical data
			\item Time-series data
		\end{enumerate}
		
		\vspace{0.1cm}
		\textcolor{myblue}{\textbf{ANSWER: B}}
		
		\vspace{0.1cm}
		\textbf{DETAILED EXPLANATION:}
		\begin{itemize}
			\item Categories with NO inherent order
			\item Example: marital status, gender, eye color
			\item Cannot be used as numeric without encoding
			\item One-hot encoding is used
		\end{itemize}
		
		\vspace{0.1cm}
		\textcolor{myred}{\textbf{WHY OTHER OPTIONS ARE WRONG:}}
		\begin{itemize}
			\item \textcolor{myred}{A:} This is ordinal data
			\item \textcolor{myred}{C:} This is numerical data
			\item \textcolor{myred}{D:} This is longitudinal data
		\end{itemize}
		
		\textcolor{myred}{\textbf{EXAM TRAP:}} Nominal = \textcolor{myblue}{no natural ordering}!
	\end{frame}
	
	\begin{frame}{Q32: Ordinal Data - Tutorial}
		\fontsize{6.5}{7.5}\selectfont
		\textbf{QUESTION: What is ordinal data?}
		
		\vspace{0.1cm}
		\textbf{OPTIONS:}
		\begin{enumerate}
			\item Categorical data with natural ordering
			\item Categorical data with no natural ordering
			\item Continuous numerical data
			\item Binary data
		\end{enumerate}
		
		\vspace{0.1cm}
		\textcolor{myblue}{\textbf{ANSWER: A}}
		
		\vspace{0.1cm}
		\textbf{DETAILED EXPLANATION:}
		\begin{itemize}
			\item Categories have natural ordering
			\item Intervals between values not uniform
			\item Examples: educational level, credit ratings, satisfaction scores
			\item Can be encoded with ordinal values
		\end{itemize}
		
		\vspace{0.1cm}
		\textcolor{myred}{\textbf{WHY OTHER OPTIONS ARE WRONG:}}
		\begin{itemize}
			\item \textcolor{myred}{B:} This is nominal data
			\item \textcolor{myred}{C:} This is numerical data
			\item \textcolor{myred}{D:} Binary is a special case
		\end{itemize}
		
		\textcolor{myred}{\textbf{EXAM TRAP:}} Ordinal = \textcolor{myblue}{natural ordering of categories}!
	\end{frame}
	
	\begin{frame}{Q33: Nominal vs Ordinal - Tutorial}
		\fontsize{6.5}{7.5}\selectfont
		\textbf{QUESTION: Which is an example of ordinal data?}
		
		\vspace{0.1cm}
		\textbf{OPTIONS:}
		\begin{enumerate}
			\item Gender (Male/Female)
			\item Marital status (Single/Married/Divorced)
			\item Credit ratings (AAA, AA, A, BBB)
			\item Country of origin
		\end{enumerate}
		
		\vspace{0.1cm}
		\textcolor{myblue}{\textbf{ANSWER: C}}
		
		\vspace{0.1cm}
		\textbf{DETAILED EXPLANATION:}
		\begin{itemize}
			\item Ordinal: categories have natural ordering
			\item Credit ratings: AAA > AA > A > BBB
			\item Ordering is meaningful
			\item Intervals between ratings are not equal
		\end{itemize}
		
		\vspace{0.1cm}
		\textcolor{myred}{\textbf{WHY OTHER OPTIONS ARE WRONG:}}
		\begin{itemize}
			\item \textcolor{myred}{A:} Gender is nominal (no order)
			\item \textcolor{myred}{B:} Marital status is nominal (no order)
			\item \textcolor{myred}{D:} Country is nominal (no order)
		\end{itemize}
		
		\textcolor{myred}{\textbf{EXAM TRAP:}} Credit ratings = \textcolor{myblue}{ordinal data}!
	\end{frame}
	
	\begin{frame}{Q34: Binary Data - Tutorial}
		\fontsize{6.5}{7.5}\selectfont
		\textbf{QUESTION: What is binary data?}
		
		\vspace{0.1cm}
		\textbf{OPTIONS:}
		\begin{enumerate}
			\item Data with two possible values
			\item Data with multiple categories
			\item Continuous data
			\item Text data
		\end{enumerate}
		
		\vspace{0.1cm}
		\textcolor{myblue}{\textbf{ANSWER: A}}
		
		\vspace{0.1cm}
		\textbf{DETAILED EXPLANATION:}
		\begin{itemize}
			\item Special case of categorical data
			\item Only two possible values (0/1, Yes/No)
			\item Example: default/no default, fraud/legitimate
			\item Can be used directly in ML models
		\end{itemize}
		
		\vspace{0.1cm}
		\textcolor{myred}{\textbf{WHY OTHER OPTIONS ARE WRONG:}}
		\begin{itemize}
			\item \textcolor{myred}{B:} This is multi-category categorical
			\item \textcolor{myred}{C:} Binary is not continuous
			\item \textcolor{myred}{D:} Binary is not text
		\end{itemize}
		
		\textcolor{myred}{\textbf{EXAM TRAP:}} Binary data = \textcolor{myblue}{two possible values}!
	\end{frame}
	
	\begin{frame}{Q35: Cross-Sectional Data - Tutorial}
		\fontsize{6.5}{7.5}\selectfont
		\textbf{QUESTION: What is cross-sectional data?}
		
		\vspace{0.1cm}
		\textbf{OPTIONS:}
		\begin{enumerate}
			\item Data spanning multiple time periods
			\item Data for a single point in time
			\item Data connected spatially
			\item Data from networks
		\end{enumerate}
		
		\vspace{0.1cm}
		\textcolor{myblue}{\textbf{ANSWER: B}}
		
		\vspace{0.1cm}
		\textbf{DETAILED EXPLANATION:}
		\begin{itemize}
			\item Snapshot at a specific moment
			\item Observations are independent
			\item Order does not matter
			\item Examples: loan portfolio characteristics, survey responses
		\end{itemize}
		
		\vspace{0.1cm}
		\textcolor{myred}{\textbf{WHY OTHER OPTIONS ARE WRONG:}}
		\begin{itemize}
			\item \textcolor{myred}{A:} This is longitudinal/time-series data
			\item \textcolor{myred}{C:} This is spatial data
			\item \textcolor{myred}{D:} This is network data
		\end{itemize}
		
		\textcolor{myred}{\textbf{EXAM TRAP:}} Cross-sectional = \textcolor{myblue}{single point in time}!
	\end{frame}
	
	\begin{frame}{Q36: Longitudinal Data - Tutorial}
		\fontsize{6.5}{7.5}\selectfont
		\textbf{QUESTION: What is longitudinal data?}
		
		\vspace{0.1cm}
		\textbf{OPTIONS:}
		\begin{enumerate}
			\item Data for a single point in time
			\item Data spanning multiple time periods or connected spatially
			\item Data that is always cross-sectional
			\item Data that cannot be analyzed
		\end{enumerate}
		
		\vspace{0.1cm}
		\textcolor{myblue}{\textbf{ANSWER: B}}
		
		\vspace{0.1cm}
		\textbf{DETAILED EXPLANATION:}
		\begin{itemize}
			\item Observations are related temporally or spatially
			\item Time-series is the most common example
			\item Order matters in analysis
			\item Examples: stock prices, economic data
		\end{itemize}
		
		\vspace{0.1cm}
		\textcolor{myred}{\textbf{WHY OTHER OPTIONS ARE WRONG:}}
		\begin{itemize}
			\item \textcolor{myred}{A:} This is cross-sectional data
			\item \textcolor{myred}{C:} Longitudinal is NOT cross-sectional
			\item \textcolor{myred}{D:} Longitudinal data can be analyzed
		\end{itemize}
		
		\textcolor{myred}{\textbf{EXAM TRAP:}} Longitudinal = \textcolor{myblue}{multiple time periods or spatial connections}!
	\end{frame}
	
	\begin{frame}{Q37: Time-Series Data - Tutorial}
		\fontsize{6.5}{7.5}\selectfont
		\textbf{QUESTION: What is time-series data?}
		
		\vspace{0.1cm}
		\textbf{OPTIONS:}
		\begin{enumerate}
			\item Data measured at a single point in time
			\item Data measured over time with chronological order
			\item Data from different locations
			\item Data with no ordering
		\end{enumerate}
		
		\vspace{0.1cm}
		\textcolor{myblue}{\textbf{ANSWER: B}}
		
		\vspace{0.1cm}
		\textbf{DETAILED EXPLANATION:}
		\begin{itemize}
			\item Observations over time
			\item Chronological order matters
			\item Observations are not independent
			\item Examples: stock prices, GDP, yields
		\end{itemize}
		
		\vspace{0.1cm}
		\textcolor{myred}{\textbf{WHY OTHER OPTIONS ARE WRONG:}}
		\begin{itemize}
			\item \textcolor{myred}{A:} This is cross-sectional data
			\item \textcolor{myred}{C:} This is spatial data
			\item \textcolor{myred}{D:} Time-series has ordering
		\end{itemize}
		
		\textcolor{myred}{\textbf{EXAM TRAP:}} Time-series = \textcolor{myblue}{data over time with chronological order}!
	\end{frame}
	
	\begin{frame}{Q38: Data Measurement Scales - Tutorial}
		\fontsize{6.5}{7.5}\selectfont
		\textbf{QUESTION: What are the four scales of data measurement?}
		
		\vspace{0.1cm}
		\textbf{OPTIONS:}
		\begin{enumerate}
			\item Nominal, Ordinal, Interval, Ratio
			\item Binary, Ordinal, Interval, Ratio
			\item Nominal, Binary, Ordinal, Continuous
			\item Categorical, Numerical, Text, Image
		\end{enumerate}
		
		\vspace{0.1cm}
		\textcolor{myblue}{\textbf{ANSWER: A}}
		
		\vspace{0.1cm}
		\textbf{DETAILED EXPLANATION:}
		\begin{itemize}
			\item \textbf{Nominal:} Categories without order
			\item \textbf{Ordinal:} Categories with order
			\item \textbf{Interval:} Equal intervals, no true zero
			\item \textbf{Ratio:} Equal intervals, true zero exists
		\end{itemize}
		
		\vspace{0.1cm}
		\textcolor{myred}{\textbf{WHY OTHER OPTIONS ARE WRONG:}}
		\begin{itemize}
			\item \textcolor{myred}{B:} Missing Nominal
			\item \textcolor{myred}{C:} Missing Interval and Ratio
			\item \textcolor{myred}{D:} These are data types, not measurement scales
		\end{itemize}
		
		\textcolor{myred}{\textbf{EXAM TRAP:}} Four scales = \textcolor{myblue}{Nominal, Ordinal, Interval, Ratio}!
	\end{frame}
	
	\begin{frame}{Q39: Interval vs Ratio Data - Tutorial}
		\fontsize{6.5}{7.5}\selectfont
		\textbf{QUESTION: What distinguishes interval from ratio data?}
		
		\vspace{0.1cm}
		\textbf{OPTIONS:}
		\begin{enumerate}
			\item Ratio data has no true zero
			\item Interval data has a true zero
			\item Ratio data has a true zero; interval does not
			\item There is no difference
		\end{enumerate}
		
		\vspace{0.1cm}
		\textcolor{myblue}{\textbf{ANSWER: C}}
		
		\vspace{0.1cm}
		\textbf{DETAILED EXPLANATION:}
		\begin{itemize}
			\item \textbf{Interval:} No true zero (e.g., temperature in \textdegree C)
			\item \textbf{Ratio:} True zero exists (e.g., height, income)
			\item Ratio: 0 means absence of the quantity
			\item Interval: 0 is arbitrary (0\textdegree C doesn't mean no temperature)
		\end{itemize}
		
		\vspace{0.1cm}
		\textcolor{myred}{\textbf{WHY OTHER OPTIONS ARE WRONG:}}
		\begin{itemize}
			\item \textcolor{myred}{A:} This is reversed
			\item \textcolor{myred}{B:} This is reversed
			\item \textcolor{myred}{D:} There is a difference (true zero)
		\end{itemize}
		
		\textcolor{myred}{\textbf{EXAM TRAP:}} Ratio = \textcolor{myblue}{true zero}; Interval = \textcolor{myblue}{no true zero}!
	\end{frame}
	
	\begin{frame}{Q40: Data Types Summary - Tutorial}
		\fontsize{6.5}{7.5}\selectfont
		\textbf{QUESTION: Which data type has a true zero point?}
		
		\vspace{0.1cm}
		\textbf{OPTIONS:}
		\begin{enumerate}
			\item Nominal
			\item Ordinal
			\item Interval
			\item Ratio
		\end{enumerate}
		
		\vspace{0.1cm}
		\textcolor{myblue}{\textbf{ANSWER: D}}
		
		\vspace{0.1cm}
		\textbf{DETAILED EXPLANATION:}
		\begin{itemize}
			\item Ratio scale: true zero exists
			\item 0 means absence of the quantity
			\item Examples: height, weight, income, distance
			\item Allows for meaningful ratios (e.g., 2x as tall)
		\end{itemize}
		
		\vspace{0.1cm}
		\textcolor{myred}{\textbf{WHY OTHER OPTIONS ARE WRONG:}}
		\begin{itemize}
			\item \textcolor{myred}{A:} Nominal has no order or zero
			\item \textcolor{myred}{B:} Ordinal has order but no true zero
			\item \textcolor{myred}{C:} Interval has no true zero
		\end{itemize}
		
		\textcolor{myred}{\textbf{EXAM TRAP:}} Ratio = \textcolor{myblue}{true zero}!
	\end{frame}
	
	% ============================================================
	% SECTION 4: DATA PRE-PROCESSING - QUESTIONS 41-60
	% ============================================================
	
	\begin{frame}{Q41: Exploratory Data Analysis (EDA) - Tutorial}
		\fontsize{6.5}{7.5}\selectfont
		\textbf{QUESTION: What is Exploratory Data Analysis (EDA)?}
		
		\vspace{0.1cm}
		\textbf{OPTIONS:}
		\begin{enumerate}
			\item Building the final model
			\item The process of collecting, cleaning, visualizing, and analyzing data before modeling
			\item Deploying the model to production
			\item Testing the model on new data
		\end{enumerate}
		
		\vspace{0.1cm}
		\textcolor{myblue}{\textbf{ANSWER: B}}
		
		\vspace{0.1cm}
		\textbf{DETAILED EXPLANATION:}
		\begin{itemize}
			\item Essential first step in ML projects
			\item Includes data collection, cleaning, visualization
			\item Understand data structure and relationships
			\item Helps select features for modeling
		\end{itemize}
		
		\vspace{0.1cm}
		\textcolor{myred}{\textbf{WHY OTHER OPTIONS ARE WRONG:}}
		\begin{itemize}
			\item \textcolor{myred}{A:} EDA is before model building
			\item \textcolor{myred}{C:} This is deployment
			\item \textcolor{myred}{D:} This is testing/evaluation
		\end{itemize}
		
		\textcolor{myred}{\textbf{EXAM TRAP:}} EDA = \textcolor{myblue}{exploring data BEFORE modeling}!
	\end{frame}
	
	\begin{frame}{Q42: Data Cleaning Importance - Tutorial}
		\fontsize{6.5}{7.5}\selectfont
		\textbf{QUESTION: What percentage of time do data analysts often spend cleaning data?}
		
		\vspace{0.1cm}
		\textbf{OPTIONS:}
		\begin{enumerate}
			\item 10\%
			\item 30\%
			\item Up to 80\%
			\item 100\%
		\end{enumerate}
		
		\vspace{0.1cm}
		\textcolor{myblue}{\textbf{ANSWER: C}}
		
		\vspace{0.1cm}
		\textbf{DETAILED EXPLANATION:}
		\begin{itemize}
			\item Data cleaning is time-consuming
			\item Often up to 80\% of project time
			\item Essential for successful ML projects
			\item "Garbage in, garbage out" (GIGO)
		\end{itemize}
		
		\vspace{0.1cm}
		\textcolor{myred}{\textbf{WHY OTHER OPTIONS ARE WRONG:}}
		\begin{itemize}
			\item \textcolor{myred}{A:} Much more than 10\%
			\item \textcolor{myred}{B:} Typical is 80\%, not 30\%
			\item \textcolor{myred}{D:} Not 100\%, but very significant
		\end{itemize}
		
		\textcolor{myred}{\textbf{EXAM TRAP:}} Data cleaning = \textcolor{myblue}{up to 80\% of project time}!
	\end{frame}
	
	\begin{frame}{Q43: GIGO Acronym - Tutorial}
		\fontsize{6.5}{7.5}\selectfont
		\textbf{QUESTION: What does GIGO stand for?}
		
		\vspace{0.1cm}
		\textbf{OPTIONS:}
		\begin{enumerate}
			\item Good Inputs, Great Outputs
			\item Garbage In, Garbage Out
			\item Great Inputs, Great Outputs
			\item Graphics In, Graphics Out
		\end{enumerate}
		
		\vspace{0.1cm}
		\textcolor{myblue}{\textbf{ANSWER: B}}
		
		\vspace{0.1cm}
		\textbf{DETAILED EXPLANATION:}
		\begin{itemize}
			\item Analysis quality depends on data quality
			\item Poor data leads to poor results
			\item Emphasizes importance of data cleaning
			\item Fundamental principle in data science
		\end{itemize}
		
		\vspace{0.1cm}
		\textcolor{myred}{\textbf{WHY OTHER OPTIONS ARE WRONG:}}
		\begin{itemize}
			\item \textcolor{myred}{A:} Incorrect expansion
			\item \textcolor{myred}{C:} Incorrect expansion
			\item \textcolor{myred}{D:} Incorrect expansion
		\end{itemize}
		
		\textcolor{myred}{\textbf{EXAM TRAP:}} GIGO = \textcolor{myblue}{Garbage In, Garbage Out}!
	\end{frame}
	
	\begin{frame}{Q44: Data Cleaning Reasons - Tutorial}
		\fontsize{6.5}{7.5}\selectfont
		\textbf{QUESTION: What is NOT a reason for data cleaning?}
		
		\vspace{0.1cm}
		\textbf{OPTIONS:}
		\begin{enumerate}
			\item Inconsistent recording
			\item Duplicate observations
			\item Outliers
			\item Increasing data size
		\end{enumerate}
		
		\vspace{0.1cm}
		\textcolor{myblue}{\textbf{ANSWER: D}}
		
		\vspace{0.1cm}
		\textbf{DETAILED EXPLANATION:}
		\begin{itemize}
			\item Data cleaning removes errors, not increases data
			\item Reasons: inconsistent recording, duplicates, outliers, missing data
			\item Goal: improve data quality
			\item Data augmentation increases data size (separate technique)
		\end{itemize}
		
		\vspace{0.1cm}
		\textcolor{myred}{\textbf{WHY OTHER OPTIONS ARE WRONG:}}
		\begin{itemize}
			\item \textcolor{myred}{A:} Inconsistent recording IS a reason for cleaning
			\item \textcolor{myred}{B:} Duplicates IS a reason for cleaning
			\item \textcolor{myred}{C:} Outliers IS a reason for cleaning
		\end{itemize}
		
		\textcolor{myred}{\textbf{EXAM TRAP:}} Data cleaning = \textcolor{myblue}{quality improvement, not size increase}!
	\end{frame}
	
	\begin{frame}{Q45: Outliers Definition - Tutorial}
		\fontsize{6.5}{7.5}\selectfont
		\textbf{QUESTION: What is an outlier in a dataset?}
		
		\vspace{0.1cm}
		\textbf{OPTIONS:}
		\begin{enumerate}
			\item The average value in the dataset
			\item Observations significantly different from the rest of the data
			\item The most common value in the dataset
			\item Data that is always correct
		\end{enumerate}
		
		\vspace{0.1cm}
		\textcolor{myblue}{\textbf{ANSWER: B}}
		
		\vspace{0.1cm}
		\textbf{DETAILED EXPLANATION:}
		\begin{itemize}
			\item Observations several standard deviations from mean
			\item May be errors or genuine rare events
			\item Can significantly affect results
			\item Treatment depends on application
		\end{itemize}
		
		\vspace{0.1cm}
		\textcolor{myred}{\textbf{WHY OTHER OPTIONS ARE WRONG:}}
		\begin{itemize}
			\item \textcolor{myred}{A:} This is the mean
			\item \textcolor{myred}{C:} This is the mode
			\item \textcolor{myred}{D:} Outliers may be errors
		\end{itemize}
		
		\textcolor{myred}{\textbf{EXAM TRAP:}} Outliers = \textcolor{myblue}{significantly different observations}!
	\end{frame}
	
	\begin{frame}{Q46: Outlier Treatment - Tutorial}
		\fontsize{6.5}{7.5}\selectfont
		\textbf{QUESTION: How should outliers be treated?}
		
		\vspace{0.1cm}
		\textbf{OPTIONS:}
		\begin{enumerate}
			\item Always remove them
			\item Always keep them
			\item Depends on the specific application and context
			\item They are never a problem
		\end{enumerate}
		
		\vspace{0.1cm}
		\textcolor{myblue}{\textbf{ANSWER: C}}
		
		\vspace{0.1cm}
		\textbf{DETAILED EXPLANATION:}
		\begin{itemize}
			\item Sometimes errors $\to$ remove
			\item Sometimes genuine rare events $\to$ keep
			\item Fraud detection: outliers may be valuable
			\item No one-size-fits-all solution
		\end{itemize}
		
		\vspace{0.1cm}
		\textcolor{myred}{\textbf{WHY OTHER OPTIONS ARE WRONG:}}
		\begin{itemize}
			\item \textcolor{myred}{A:} Not always remove
			\item \textcolor{myred}{B:} Not always keep
			\item \textcolor{myred}{D:} Outliers ARE a problem to handle
		\end{itemize}
		
		\textcolor{myred}{\textbf{EXAM TRAP:}} Outlier treatment = \textcolor{myblue}{depends on application}!
	\end{frame}
	
	\begin{frame}{Q47: Outlier Robust Models - Tutorial}
		\fontsize{6.5}{7.5}\selectfont
		\textbf{QUESTION: Which models are robust to outliers?}
		
		\vspace{0.1cm}
		\textbf{OPTIONS:}
		\begin{enumerate}
			\item Linear regression
			\item Tree-based models (decision trees, random forests)
			\item Logistic regression
			\item All models are equally sensitive
		\end{enumerate}
		
		\vspace{0.1cm}
		\textcolor{myblue}{\textbf{ANSWER: B}}
		
		\vspace{0.1cm}
		\textbf{DETAILED EXPLANATION:}
		\begin{itemize}
			\item Tree-based models split data into regions
			\item Outliers get their own branches
			\item Less affected by extreme values
			\item Chapter 4 covers these models
		\end{itemize}
		
		\vspace{0.1cm}
		\textcolor{myred}{\textbf{WHY OTHER OPTIONS ARE WRONG:}}
		\begin{itemize}
			\item \textcolor{myred}{A:} Linear regression is sensitive to outliers
			\item \textcolor{myred}{C:} Logistic regression can be sensitive
			\item \textcolor{myred}{D:} Models have different sensitivities
		\end{itemize}
		
		\textcolor{myred}{\textbf{EXAM TRAP:}} Tree-based models = \textcolor{myblue}{robust to outliers}!
	\end{frame}
	
	\begin{frame}{Q48: Missing Data Imputation - Tutorial}
		\fontsize{6.5}{7.5}\selectfont
		\textbf{QUESTION: What is imputation in data cleaning?}
		
		\vspace{0.1cm}
		\textbf{OPTIONS:}
		\begin{enumerate}
			\item Removing all rows with missing data
			\item Replacing missing values with estimated values
			\item Ignoring missing values
			\item Creating more missing values
		\end{enumerate}
		
		\vspace{0.1cm}
		\textcolor{myblue}{\textbf{ANSWER: B}}
		
		\vspace{0.1cm}
		\textbf{DETAILED EXPLANATION:}
		\begin{itemize}
			\item Replace missing observations with estimates
			\item Common methods: mean, median, mode imputation
			\item More sophisticated: model-based imputation
			\item Alternative to deleting observations
		\end{itemize}
		
		\vspace{0.1cm}
		\textcolor{myred}{\textbf{WHY OTHER OPTIONS ARE WRONG:}}
		\begin{itemize}
			\item \textcolor{myred}{A:} This is deletion, not imputation
			\item \textcolor{myred}{C:} Ignoring is not a solution
			\item \textcolor{myred}{D:} Imputation FILLS missing values
		\end{itemize}
		
		\textcolor{myred}{\textbf{EXAM TRAP:}} Imputation = \textcolor{myblue}{replacing missing values with estimates}!
	\end{frame}
	
	\begin{frame}{Q49: Informative Missingness - Tutorial}
		\fontsize{6.5}{7.5}\selectfont
		\textbf{QUESTION: What is informative missingness?}
		
		\vspace{0.1cm}
		\textbf{OPTIONS:}
		\begin{enumerate}
			\item Missing data that is random
			\item Missing data that is associated with the outcome
			\item Missing data that is always structural
			\item Missing data that cannot be imputed
		\end{enumerate}
		
		\vspace{0.1cm}
		\textcolor{myblue}{\textbf{ANSWER: B}}
		
		\vspace{0.1cm}
		\textbf{DETAILED EXPLANATION:}
		\begin{itemize}
			\item Pattern of missing data relates to outcome
			\item Can induce significant bias
			\item Example: credit card data only for solvent clients
			\item Requires careful handling
		\end{itemize}
		
		\vspace{0.1cm}
		\textcolor{myred}{\textbf{WHY OTHER OPTIONS ARE WRONG:}}
		\begin{itemize}
			\item \textcolor{myred}{A:} Random missingness is not informative
			\item \textcolor{myred}{C:} Structural missing is different
			\item \textcolor{myred}{D:} It can be handled carefully
		\end{itemize}
		
		\textcolor{myred}{\textbf{EXAM TRAP:}} Informative missingness = \textcolor{myblue}{missing pattern related to outcome}!
	\end{frame}
	
	\begin{frame}{Q50: Data Visualization Tools - Tutorial}
		\fontsize{6.5}{7.5}\selectfont
		\textbf{QUESTION: Which is a data visualization tool?}
		
		\vspace{0.1cm}
		\textbf{OPTIONS:}
		\begin{enumerate}
			\item Histogram
			\item Correlation matrix
			\item Scatter plot
			\item All of the above
		\end{enumerate}
		
		\vspace{0.1cm}
		\textcolor{myblue}{\textbf{ANSWER: D}}
		
		\vspace{0.1cm}
		\textbf{DETAILED EXPLANATION:}
		\begin{itemize}
			\item \textbf{Histogram:} Distribution of a single variable
			\item \textbf{Scatter plot:} Relationship between two variables
			\item \textbf{Correlation matrix:} Relationships among multiple variables
			\item \textbf{Boxplot:} Summary of distribution with outliers
		\end{itemize}
		
		\vspace{0.1cm}
		\textcolor{myred}{\textbf{WHY OTHER OPTIONS ARE WRONG:}}
		\begin{itemize}
			\item \textcolor{myred}{A:} Correct but incomplete
			\item \textcolor{myred}{B:} Correct but incomplete
			\item \textcolor{myred}{C:} Correct but incomplete
		\end{itemize}
		
		\textcolor{myred}{\textbf{EXAM TRAP:}} All are \textcolor{myblue}{data visualization tools}!
	\end{frame}
	
	\begin{frame}{Q51: Histogram Purpose - Tutorial}
		\fontsize{6.5}{7.5}\selectfont
		\textbf{QUESTION: What does a histogram show?}
		
		\vspace{0.1cm}
		\textbf{OPTIONS:}
		\begin{enumerate}
			\item Relationship between two variables
			\item The distribution/shape of a single variable
			\item The correlation between features
			\item The missing values in the data
		\end{enumerate}
		
		\vspace{0.1cm}
		\textcolor{myblue}{\textbf{ANSWER: B}}
		
		\vspace{0.1cm}
		\textbf{DETAILED EXPLANATION:}
		\begin{itemize}
			\item Count of observations in bins
			\item Shows distribution shape
			\item Identifies skewness, outliers, patterns
			\item Area under histogram equals 1 (when density)
		\end{itemize}
		
		\vspace{0.1cm}
		\textcolor{myred}{\textbf{WHY OTHER OPTIONS ARE WRONG:}}
		\begin{itemize}
			\item \textcolor{myred}{A:} This is scatter plot
			\item \textcolor{myred}{C:} This is correlation matrix
			\item \textcolor{myred}{D:} This is not what histogram shows
		\end{itemize}
		
		\textcolor{myred}{\textbf{EXAM TRAP:}} Histogram shows \textcolor{myblue}{distribution of one variable}!
	\end{frame}
	
	\begin{frame}{Q52: Skewness Definition - Tutorial}
		\fontsize{6.5}{7.5}\selectfont
		\textbf{QUESTION: What is skewness in a distribution?}
		
		\vspace{0.1cm}
		\textbf{OPTIONS:}
		\begin{enumerate}
			\item Symmetric distribution
			\item Asymmetric distribution (longer tail on one side)
			\item Distribution with no tails
			\item Distribution with multiple peaks
		\end{enumerate}
		
		\vspace{0.1cm}
		\textcolor{myblue}{\textbf{ANSWER: B}}
		
		\vspace{0.1cm}
		\textbf{DETAILED EXPLANATION:}
		\begin{itemize}
			\item \textbf{Negative skew:} Longer left tail
			\item \textbf{Positive skew:} Longer right tail
			\item Example: stock returns have negative skew
			\item Can be addressed with transformations (log)
		\end{itemize}
		
		\vspace{0.1cm}
		\textcolor{myred}{\textbf{WHY OTHER OPTIONS ARE WRONG:}}
		\begin{itemize}
			\item \textcolor{myred}{A:} Symmetric is no skew
			\item \textcolor{myred}{C:} All distributions have tails
			\item \textcolor{myred}{D:} This is multimodal
		\end{itemize}
		
		\textcolor{myred}{\textbf{EXAM TRAP:}} Skewness = \textcolor{myblue}{asymmetric distribution}!
	\end{frame}
	
	\begin{frame}{Q53: Q-Q Plot Purpose - Tutorial}
		\fontsize{6.5}{7.5}\selectfont
		\textbf{QUESTION: What is a Q-Q plot used for?}
		
		\vspace{0.1cm}
		\textbf{OPTIONS:}
		\begin{enumerate}
			\item Showing correlations
			\item Comparing empirical distribution to theoretical distribution
			\item Showing missing values
			\item Showing outliers only
		\end{enumerate}
		
		\vspace{0.1cm}
		\textcolor{myblue}{\textbf{ANSWER: B}}
		
		\vspace{0.1cm}
		\textbf{DETAILED EXPLANATION:}
		\begin{itemize}
			\item Q-Q = Quantile-Quantile
			\item Compares empirical vs theoretical quantiles
			\item Points on identity line = distribution matches
			\item Normal Q-Q plot: checks normality
		\end{itemize}
		
		\vspace{0.1cm}
		\textcolor{myred}{\textbf{WHY OTHER OPTIONS ARE WRONG:}}
		\begin{itemize}
			\item \textcolor{myred}{A:} Correlation is different
			\item \textcolor{myred}{C:} Missing values are different
			\item \textcolor{myred}{D:} Shows more than just outliers
		\end{itemize}
		
		\textcolor{myred}{\textbf{EXAM TRAP:}} Q-Q plot = \textcolor{myblue}{empirical vs theoretical distribution}!
	\end{frame}
	
	\begin{frame}{Q54: Boxplot Purpose - Tutorial}
		\fontsize{6.5}{7.5}\selectfont
		\textbf{QUESTION: What does a boxplot show?}
		
		\vspace{0.1cm}
		\textbf{OPTIONS:}
		\begin{enumerate}
			\item Mean and standard deviation only
			\item Five-number summary and outliers
			\item Correlation between variables
			\item The relationship between two variables
		\end{enumerate}
		
		\vspace{0.1cm}
		\textcolor{myblue}{\textbf{ANSWER: B}}
		
		\vspace{0.1cm}
		\textbf{DETAILED EXPLANATION:}
		\begin{itemize}
			\item Box = 1st quartile to 3rd quartile
			\item Line inside = median
			\item Whiskers = 1.5 $\times$ IQR
			\item Points beyond whiskers = outliers
		\end{itemize}
		
		\vspace{0.1cm}
		\textcolor{myred}{\textbf{WHY OTHER OPTIONS ARE WRONG:}}
		\begin{itemize}
			\item \textcolor{myred}{A:} Boxplot shows quartiles, not just mean/sd
			\item \textcolor{myred}{C:} Correlation is different
			\item \textcolor{myred}{D:} Scatter plot shows relationships
		\end{itemize}
		
		\textcolor{myred}{\textbf{EXAM TRAP:}} Boxplot = \textcolor{myblue}{five-number summary + outliers}!
	\end{frame}
	
	\begin{frame}{Q55: Correlation Matrix - Tutorial}
		\fontsize{6.5}{7.5}\selectfont
		\textbf{QUESTION: What does a correlation matrix show?}
		
		\vspace{0.1cm}
		\textbf{OPTIONS:}
		\begin{enumerate}
			\item Distribution of each variable
			\item Pairwise correlations between variables
			\item Missing values
			\item Outliers
		\end{enumerate}
		
		\vspace{0.1cm}
		\textcolor{myblue}{\textbf{ANSWER: B}}
		
		\vspace{0.1cm}
		\textbf{DETAILED EXPLANATION:}
		\begin{itemize}
			\item Matrix of pairwise correlations
			\item Values between -1 and 1
			\item Diagonal = 1 (variable with itself)
			\item Helps identify highly correlated features
		\end{itemize}
		
		\vspace{0.1cm}
		\textcolor{myred}{\textbf{WHY OTHER OPTIONS ARE WRONG:}}
		\begin{itemize}
			\item \textcolor{myred}{A:} Histogram shows distribution
			\item \textcolor{myred}{C:} Missing values are different
			\item \textcolor{myred}{D:} Outliers are different
		\end{itemize}
		
		\textcolor{myred}{\textbf{EXAM TRAP:}} Correlation matrix = \textcolor{myblue}{pairwise correlations}!
	\end{frame}
	
	\begin{frame}{Q56: Standardization Formula - Tutorial}
		\fontsize{6.5}{7.5}\selectfont
		\textbf{QUESTION: What is the formula for standardization?}
		
		\vspace{0.1cm}
		\textbf{OPTIONS:}
		\begin{enumerate}
			\item $\tilde{x}_{ij} = \frac{x_{ij} - \min(x_i)}{\max(x_i) - \min(x_i)}$
			\item $\tilde{x}_{ij} = \frac{x_{ij} - \hat{\mu}_i}{\hat{\sigma}_i}$
			\item $\tilde{x}_{ij} = x_{ij} - \hat{\mu}_i$
			\item $\tilde{x}_{ij} = \frac{x_{ij}}{\hat{\sigma}_i}$
		\end{enumerate}
		
		\vspace{0.1cm}
		\textcolor{myblue}{\textbf{ANSWER: B}}
		
		\vspace{0.1cm}
		\textbf{DETAILED EXPLANATION:}
		\begin{itemize}
			\item $\hat{\mu}_i$: sample mean of variable i
			\item $\hat{\sigma}_i$: sample standard deviation
			\item New variable: mean = 0, variance = 1
			\item Can take any value (positive/negative)
		\end{itemize}
		
		\vspace{0.1cm}
		\textcolor{myred}{\textbf{WHY OTHER OPTIONS ARE WRONG:}}
		\begin{itemize}
			\item \textcolor{myred}{A:} This is normalization (min-max)
			\item \textcolor{myred}{C:} Only subtracts mean (centering)
			\item \textcolor{myred}{D:} Only divides by standard deviation
		\end{itemize}
		
		\textcolor{myred}{\textbf{EXAM TRAP:}} Standardization = \textcolor{myblue}{$(x - mean) / sd$}!
	\end{frame}
	
	\begin{frame}{Q57: Normalization Formula - Tutorial}
		\fontsize{6.5}{7.5}\selectfont
		\textbf{QUESTION: What is the formula for normalization (min-max transformation)?}
		
		\vspace{0.1cm}
		\textbf{OPTIONS:}
		\begin{enumerate}
			\item $\tilde{x}_{ij} = \frac{x_{ij} - \min(x_i)}{\max(x_i) - \min(x_i)}$
			\item $\tilde{x}_{ij} = \frac{x_{ij} - \hat{\mu}_i}{\hat{\sigma}_i}$
			\item $\tilde{x}_{ij} = \frac{x_{ij}}{\hat{\mu}_i}$
			\item $\tilde{x}_{ij} = x_{ij} - \hat{\mu}_i$
		\end{enumerate}
		
		\vspace{0.1cm}
		\textcolor{myblue}{\textbf{ANSWER: A}}
		
		\vspace{0.1cm}
		\textbf{DETAILED EXPLANATION:}
		\begin{itemize}
			\item Bounds between 0 and 1
			\item Uses min and max of the variable
			\item All values in [0, 1] range
			\item Mean not necessarily 0, variance not necessarily 1
		\end{itemize}
		
		\vspace{0.1cm}
		\textcolor{myred}{\textbf{WHY OTHER OPTIONS ARE WRONG:}}
		\begin{itemize}
			\item \textcolor{myred}{B:} This is standardization
			\item \textcolor{myred}{C:} Not a standard method
			\item \textcolor{myred}{D:} This is centering only
		\end{itemize}
		
		\textcolor{myred}{\textbf{EXAM TRAP:}} Normalization = \textcolor{myblue}{$(x - min) / (max - min)$}!
	\end{frame}
	
	\begin{frame}{Q58: Standardization vs Normalization - Tutorial}
		\fontsize{6.5}{7.5}\selectfont
		\textbf{QUESTION: When is standardization preferred over normalization?}
		
		\vspace{0.1cm}
		\textbf{OPTIONS:}
		\begin{enumerate}
			\item When data has no outliers
			\item When data contains outliers
			\item When data is categorical
			\item When data is already in [0,1] range
		\end{enumerate}
		
		\vspace{0.1cm}
		\textcolor{myblue}{\textbf{ANSWER: B}}
		
		\vspace{0.1cm}
		\textbf{DETAILED EXPLANATION:}
		\begin{itemize}
			\item Standardization handles outliers better
			\item Normalization squashes data into tight range
			\item Outliers can distort min-max range
			\item Standardization uses mean and sd (more robust)
		\end{itemize}
		
		\vspace{0.1cm}
		\textcolor{myred}{\textbf{WHY OTHER OPTIONS ARE WRONG:}}
		\begin{itemize}
			\item \textcolor{myred}{A:} Normalization works without outliers
			\item \textcolor{myred}{C:} Both are for numerical data
			\item \textcolor{myred}{D:} Scaling not needed if already [0,1]
		\end{itemize}
		
		\textcolor{myred}{\textbf{EXAM TRAP:}} Standardization = \textcolor{myblue}{better for data with outliers}!
	\end{frame}
	
	\begin{frame}{Q59: Reasons for Scaling - Tutorial}
		\fontsize{6.5}{7.5}\selectfont
		\textbf{QUESTION: What is NOT a reason for scaling data?}
		
		\vspace{0.1cm}
		\textbf{OPTIONS:}
		\begin{enumerate}
			\item Numerical stability of learning algorithm
			\item Ease of interpretation of parameter estimates
			\item To make data categorical
			\item To check if predictions are extrapolations
		\end{enumerate}
		
		\vspace{0.1cm}
		\textcolor{myblue}{\textbf{ANSWER: C}}
		
		\vspace{0.1cm}
		\textbf{DETAILED EXPLANATION:}
		\begin{itemize}
			\item Scaling doesn't make data categorical
			\item Reasons: numerical stability, interpretation, extrapolation detection
			\item Scaling transforms numerical values
			\item Categorical variables need encoding (not scaling)
		\end{itemize}
		
		\vspace{0.1cm}
		\textcolor{myred}{\textbf{WHY OTHER OPTIONS ARE WRONG:}}
		\begin{itemize}
			\item \textcolor{myred}{A:} This IS a reason for scaling
			\item \textcolor{myred}{B:} This IS a reason for scaling
			\item \textcolor{myred}{D:} This IS a reason for scaling
		\end{itemize}
		
		\textcolor{myred}{\textbf{EXAM TRAP:}} Scaling = \textcolor{myblue}{numerical transformation, not categorization}!
	\end{frame}
	
	\begin{frame}{Q60: Log Transformation - Tutorial}
		\fontsize{6.5}{7.5}\selectfont
		\textbf{QUESTION: Why apply a log transformation to skewed data?}
		
		\vspace{0.1cm}
		\textbf{OPTIONS:}
		\begin{enumerate}
			\item To make data categorical
			\item To reduce skewness and make distribution more symmetric
			\item To increase skewness
			\item To remove all outliers
		\end{enumerate}
		
		\vspace{0.1cm}
		\textcolor{myblue}{\textbf{ANSWER: B}}
		
		\vspace{0.1cm}
		\textbf{DETAILED EXPLANATION:}
		\begin{itemize}
			\item Log transformation reduces positive skew
			\item Data becomes better behaved
			\item Example: trading volume distribution becomes more symmetric
			\item Rule of thumb: if max/min $>$ 10, consider log transform
		\end{itemize}
		
		\vspace{0.1cm}
		\textcolor{myred}{\textbf{WHY OTHER OPTIONS ARE WRONG:}}
		\begin{itemize}
			\item \textcolor{myred}{A:} Log doesn't make data categorical
			\item \textcolor{myred}{C:} Log REDUCES skewness
			\item \textcolor{myred}{D:} Log doesn't remove outliers
		\end{itemize}
		
		\textcolor{myred}{\textbf{EXAM TRAP:}} Log transform = \textcolor{myblue}{reduce skewness}!
	\end{frame}
	
	% ============================================================
	% SECTION 5: FEATURE EXTRACTION - QUESTIONS 61-70
	% ============================================================
	
	\begin{frame}{Q61: Feature Extraction - Tutorial}
		\fontsize{6.5}{7.5}\selectfont
		\textbf{QUESTION: What is feature extraction/engineering?}
		
		\vspace{0.1cm}
		\textbf{OPTIONS:}
		\begin{enumerate}
			\item Removing all features
			\item Transforming raw attributes into features suitable for analysis
			\item Always using raw data without modification
			\item Only using categorical features
		\end{enumerate}
		
		\vspace{0.1cm}
		\textcolor{myblue}{\textbf{ANSWER: B}}
		
		\vspace{0.1cm}
		\textbf{DETAILED EXPLANATION:}
		\begin{itemize}
			\item Process of creating useful features from raw data
			\item Includes encoding categorical variables
			\item May include creating derived features
			\item Critical step in ML pipeline
		\end{itemize}
		
		\vspace{0.1cm}
		\textcolor{myred}{\textbf{WHY OTHER OPTIONS ARE WRONG:}}
		\begin{itemize}
			\item \textcolor{myred}{A:} Feature extraction creates, not removes
			\item \textcolor{myred}{C:} Raw data often needs transformation
			\item \textcolor{myred}{D:} Both numerical and categorical features used
		\end{itemize}
		
		\textcolor{myred}{\textbf{EXAM TRAP:}} Feature extraction = \textcolor{myblue}{transforming raw attributes into features}!
	\end{frame}
	
	\begin{frame}{Q62: One-Hot Encoding - Tutorial}
		\fontsize{6.5}{7.5}\selectfont
		\textbf{QUESTION: What is one-hot encoding?}
		
		\vspace{0.1cm}
		\textbf{OPTIONS:}
		\begin{enumerate}
			\item Converting numerical data to categorical
			\item Creating binary dummy variables for each category
			\item Removing categorical variables
			\item Converting text to numbers
		\end{enumerate}
		
		\vspace{0.1cm}
		\textcolor{myblue}{\textbf{ANSWER: B}}
		
		\vspace{0.1cm}
		\textbf{DETAILED EXPLANATION:}
		\begin{itemize}
			\item Also called binarization
			\item Each category becomes a binary (0/1) variable
			\item Only one dummy = 1 for each observation
			\item Essential for nominal categorical variables
		\end{itemize}
		
		\vspace{0.1cm}
		\textcolor{myred}{\textbf{WHY OTHER OPTIONS ARE WRONG:}}
		\begin{itemize}
			\item \textcolor{myred}{A:} One-hot encodes categorical to binary
			\item \textcolor{myred}{C:} It encodes, not removes
			\item \textcolor{myred}{D:} Text encoding is different (NLP)
		\end{itemize}
		
		\textcolor{myred}{\textbf{EXAM TRAP:}} One-hot encoding = \textcolor{myblue}{binary dummies for each category}!
	\end{frame}
	
	\begin{frame}{Q63: One-Hot Encoding Example - Tutorial}
		\fontsize{6.5}{7.5}\selectfont
		\textbf{QUESTION: How would you encode "Marital Status" (Single/Married/Divorced)?}
		
		\vspace{0.1cm}
		\textbf{OPTIONS:}
		\begin{enumerate}
			\item Single=0, Married=1, Divorced=2
			\item Create three binary dummies: Single, Married, Divorced
			\item Remove the variable
			\item Convert to text
		\end{enumerate}
		
		\vspace{0.1cm}
		\textcolor{myblue}{\textbf{ANSWER: B}}
		
		\vspace{0.1cm}
		\textbf{DETAILED EXPLANATION:}
		\begin{itemize}
			\item Marital status is nominal (no ordering)
			\item One-hot: three binary variables
			\item Single: [1,0,0], Married: [0,1,0], Divorced: [0,0,1]
			\item Avoids implying ordering between categories
		\end{itemize}
		
		\vspace{0.1cm}
		\textcolor{myred}{\textbf{WHY OTHER OPTIONS ARE WRONG:}}
		\begin{itemize}
			\item \textcolor{myred}{A:} This implies ordering (incorrect for nominal)
			\item \textcolor{myred}{C:} Removing loses valuable information
			\item \textcolor{myred}{D:} This doesn't create usable features
		\end{itemize}
		
		\textcolor{myred}{\textbf{EXAM TRAP:}} Nominal data = \textcolor{myblue}{one-hot encoding} (not numeric assignment)!
	\end{frame}
	
	\begin{frame}{Q64: Dummy Variable Trap - Tutorial}
		\fontsize{6.5}{7.5}\selectfont
		\textbf{QUESTION: What is the dummy variable trap?}
		
		\vspace{0.1cm}
		\textbf{OPTIONS:}
		\begin{enumerate}
			\item Using too many dummy variables
			\item Creating perfect multicollinearity when including all dummies and intercept
			\item Not including enough dummy variables
			\item Using only one dummy variable
		\end{enumerate}
		
		\vspace{0.1cm}
		\textcolor{myblue}{\textbf{ANSWER: B}}
		
		\vspace{0.1cm}
		\textbf{DETAILED EXPLANATION:}
		\begin{itemize}
			\item Including all dummies + intercept creates perfect collinearity
			\item Example: Male + Female + Intercept = 1 for all observations
			\item Solution: omit one dummy or set intercept to zero
			\item Also known as the reference category approach
		\end{itemize}
		
		\vspace{0.1cm}
		\textcolor{myred}{\textbf{WHY OTHER OPTIONS ARE WRONG:}}
		\begin{itemize}
			\item \textcolor{myred}{A:} Too many is not the trap
			\item \textcolor{myred}{C:} Not enough is not the trap
			\item \textcolor{myred}{D:} One dummy is fine
		\end{itemize}
		
		\textcolor{myred}{\textbf{EXAM TRAP:}} Dummy variable trap = \textcolor{myblue}{perfect multicollinearity with intercept}!
	\end{frame}
	
	\begin{frame}{Q65: Dummy Variable Trap Solutions - Tutorial}
		\fontsize{6.5}{7.5}\selectfont
		\textbf{QUESTION: How can the dummy variable trap be avoided?}
		
		\vspace{0.1cm}
		\textbf{OPTIONS:}
		\begin{enumerate}
			\item Include all dummies and intercept
			\item Omit one dummy variable or set intercept to zero
			\item Always use two dummies only
			\item Never use dummy variables
		\end{enumerate}
		
		\vspace{0.1cm}
		\textcolor{myblue}{\textbf{ANSWER: B}}
		
		\vspace{0.1cm}
		\textbf{DETAILED EXPLANATION:}
		\begin{itemize}
			\item \textbf{Omit one dummy:} Reference category
			\item \textbf{Set intercept to zero:} Include all dummies
			\item Both solutions avoid perfect multicollinearity
			\item Most common: omit one dummy variable
		\end{itemize}
		
		\vspace{0.1cm}
		\textcolor{myred}{\textbf{WHY OTHER OPTIONS ARE WRONG:}}
		\begin{itemize}
			\item \textcolor{myred}{A:} This creates the trap
			\item \textcolor{myred}{C:} May not be sufficient
			\item \textcolor{myred}{D:} Dummies are often needed
		\end{itemize}
		
		\textcolor{myred}{\textbf{EXAM TRAP:}} Solutions = \textcolor{myblue}{omit one dummy or set intercept to zero}!
	\end{frame}
	
	\begin{frame}{Q66: Ordinal Encoding - Tutorial}
		\fontsize{6.5}{7.5}\selectfont
		\textbf{QUESTION: How should ordinal categorical data be encoded?}
		
		\vspace{0.1cm}
		\textbf{OPTIONS:}
		\begin{enumerate}
			\item Always use one-hot encoding
			\item Use numeric values that preserve the ordering
			\item Remove the variable
			\item Convert to text
		\end{enumerate}
		
		\vspace{0.1cm}
		\textcolor{myblue}{\textbf{ANSWER: B}}
		
		\vspace{0.1cm}
		\textbf{DETAILED EXPLANATION:}
		\begin{itemize}
			\item Ordinal has natural ordering
			\item Assign values: 0, 1, 2, ... preserving order
			\item Example: Education (0=HS, 1=College, 2=Post-grad)
			\item Can be used as feature in models
		\end{itemize}
		
		\vspace{0.1cm}
		\textcolor{myred}{\textbf{WHY OTHER OPTIONS ARE WRONG:}}
		\begin{itemize}
			\item \textcolor{myred}{A:} One-hot loses ordering information
			\item \textcolor{myred}{C:} Removing loses valuable information
			\item \textcolor{myred}{D:} Text is not usable directly
		\end{itemize}
		
		\textcolor{myred}{\textbf{EXAM TRAP:}} Ordinal = \textcolor{myblue}{numeric encoding preserving order}!
	\end{frame}
	
	\begin{frame}{Q67: Discretization Definition - Tutorial}
		\fontsize{6.5}{7.5}\selectfont
		\textbf{QUESTION: What is discretization (binning)?}
		
		\vspace{0.1cm}
		\textbf{OPTIONS:}
		\begin{enumerate}
			\item Converting continuous data to categorical data by grouping into ranges
			\item Converting categorical to numerical
			\item Removing outliers
			\item Increasing data size
		\end{enumerate}
		
		\vspace{0.1cm}
		\textcolor{myblue}{\textbf{ANSWER: A}}
		
		\vspace{0.1cm}
		\textbf{DETAILED EXPLANATION:}
		\begin{itemize}
			\item Also called binning
			\item Group continuous values into categories
			\item Example: Firm size by market cap ranges
			\item Useful when exact values don't matter
		\end{itemize}
		
		\vspace{0.1cm}
		\textcolor{myred}{\textbf{WHY OTHER OPTIONS ARE WRONG:}}
		\begin{itemize}
			\item \textcolor{myred}{B:} This is encoding, not discretization
			\item \textcolor{myred}{C:} Outliers are different
			\item \textcolor{myred}{D:} Discretization doesn't increase size
		\end{itemize}
		
		\textcolor{myred}{\textbf{EXAM TRAP:}} Discretization = \textcolor{myblue}{continuous $\to$ categorical via binning}!
	\end{frame}
	
	\begin{frame}{Q68: Discretization Example - Tutorial}
		\fontsize{6.5}{7.5}\selectfont
		\textbf{QUESTION: When would you discretize a continuous variable?}
		
		\vspace{0.1cm}
		\textbf{OPTIONS:}
		\begin{enumerate}
			\item When exact values are important
			\item When you only care about ranges (e.g., capital above/below threshold)
			\item Always discretize all variables
			\item Never discretize continuous variables
		\end{enumerate}
		
		\vspace{0.1cm}
		\textcolor{myblue}{\textbf{ANSWER: B}}
		
		\vspace{0.1cm}
		\textbf{DETAILED EXPLANATION:}
		\begin{itemize}
			\item Example: Regulatory capital - above/below threshold
			\item Or market cap ranges (Micro, Small, Medium, Large)
			\item Useful when categories are meaningful
			\item Loss of information if exact values matter
		\end{itemize}
		
		\vspace{0.1cm}
		\textcolor{myred}{\textbf{WHY OTHER OPTIONS ARE WRONG:}}
		\begin{itemize}
			\item \textcolor{myred}{A:} If exact values matter, keep continuous
			\item \textcolor{myred}{C:} Should only discretize when meaningful
			\item \textcolor{myred}{D:} Sometimes discretization is useful
		\end{itemize}
		
		\textcolor{myred}{\textbf{EXAM TRAP:}} Discretize when \textcolor{myblue}{categories are more meaningful than exact values}!
	\end{frame}
	
	\begin{frame}{Q69: Feature Extraction Summary - Tutorial}
		\fontsize{6.5}{7.5}\selectfont
		\textbf{QUESTION: Which is NOT a method of feature extraction?}
		
		\vspace{0.1cm}
		\textbf{OPTIONS:}
		\begin{enumerate}
			\item One-hot encoding categorical variables
			\item Applying a log transformation
			\item Discretizing continuous variables
			\item Deleting all features
		\end{enumerate}
		
		\vspace{0.1cm}
		\textcolor{myblue}{\textbf{ANSWER: D}}
		
		\vspace{0.1cm}
		\textbf{DETAILED EXPLANATION:}
		\begin{itemize}
			\item Feature extraction creates/transforms features
			\item One-hot encoding: categorical to binary
			\item Log transformation: numerical to reduce skew
			\item Discretization: continuous to categorical
			\item Deleting features is feature selection, not extraction
		\end{itemize}
		
		\vspace{0.1cm}
		\textcolor{myred}{\textbf{WHY OTHER OPTIONS ARE WRONG:}}
		\begin{itemize}
			\item \textcolor{myred}{A:} This IS feature extraction
			\item \textcolor{myred}{B:} This IS feature extraction
			\item \textcolor{myred}{C:} This IS feature extraction
		\end{itemize}
		
		\textcolor{myred}{\textbf{EXAM TRAP:}} Feature extraction = \textcolor{myblue}{creating/transforming, not deleting}!
	\end{frame}
	
	\begin{frame}{Q70: Encoding Methods Table - Tutorial}
		\fontsize{6.5}{7.5}\selectfont
		\textbf{QUESTION: Which encoding method is appropriate for nominal data?}
		
		\vspace{0.1cm}
		\textbf{OPTIONS:}
		\begin{enumerate}
			\item Ordinal encoding (0,1,2,...)
			\item One-hot encoding (binary dummies)
			\item Log transformation
			\item Standardization
		\end{enumerate}
		
		\vspace{0.1cm}
		\textcolor{myblue}{\textbf{ANSWER: B}}
		
		\vspace{0.1cm}
		\textbf{DETAILED EXPLANATION:}
		\begin{itemize}
			\item Nominal data has no natural ordering
			\item One-hot encoding: separate binary variable for each category
			\item Avoids implying false ordering
			\item Examples: marital status, gender, region
		\end{itemize}
		
		\vspace{0.1cm}
		\textcolor{myred}{\textbf{WHY OTHER OPTIONS ARE WRONG:}}
		\begin{itemize}
			\item \textcolor{myred}{A:} Ordinal encoding assumes ordering (wrong for nominal)
			\item \textcolor{myred}{C:} Log is for numerical data
			\item \textcolor{myred}{D:} Standardization is for numerical data
		\end{itemize}
		
		\textcolor{myred}{\textbf{EXAM TRAP:}} Nominal data = \textcolor{myblue}{one-hot encoding}!
	\end{frame}
	
	% ============================================================
	% SECTION 6: PCA - QUESTIONS 71-85
	% ============================================================
	
	\begin{frame}{Q71: PCA Definition - Tutorial}
		\fontsize{6.5}{7.5}\selectfont
		\textbf{QUESTION: What is Principal Components Analysis (PCA)?}
		
		\vspace{0.1cm}
		\textbf{OPTIONS:}
		\begin{enumerate}
			\item A supervised learning method
			\item A dimensionality reduction technique that finds linear combinations of original features
			\item A clustering algorithm
			\item A method for handling missing data
		\end{enumerate}
		
		\vspace{0.1cm}
		\textcolor{myblue}{\textbf{ANSWER: B}}
		
		\vspace{0.1cm}
		\textbf{DETAILED EXPLANATION:}
		\begin{itemize}
			\item Most common dimensionality reduction technique
			\item Finds linear combinations of original predictors
			\item Summarizes variability in data
			\item Principal components are uncorrelated
		\end{itemize}
		
		\vspace{0.1cm}
		\textcolor{myred}{\textbf{WHY OTHER OPTIONS ARE WRONG:}}
		\begin{itemize}
			\item \textcolor{myred}{A:} PCA is UNSUPERVISED learning
			\item \textcolor{myred}{C:} PCA is for dimensionality reduction, not clustering
			\item \textcolor{myred}{D:} Missing data handling is different
		\end{itemize}
		
		\textcolor{myred}{\textbf{EXAM TRAP:}} PCA = \textcolor{myblue}{unsupervised dimensionality reduction}!
	\end{frame}
	
	\begin{frame}{Q72: PCA Purpose - Tutorial}
		\fontsize{6.5}{7.5}\selectfont
		\textbf{QUESTION: What is the primary purpose of PCA?}
		
		\vspace{0.1cm}
		\textbf{OPTIONS:}
		\begin{enumerate}
			\item To increase the number of features
			\item To reduce dimensionality while preserving variability
			\item To classify observations
			\item To cluster data
		\end{enumerate}
		
		\vspace{0.1cm}
		\textcolor{myblue}{\textbf{ANSWER: B}}
		
		\vspace{0.1cm}
		\textbf{DETAILED EXPLANATION:}
		\begin{itemize}
			\item Reduce number of features (dimensions)
			\item Preserve most variability in data
			\item Components ordered by explained variance
			\item First PC explains most variance
		\end{itemize}
		
		\vspace{0.1cm}
		\textcolor{myred}{\textbf{WHY OTHER OPTIONS ARE WRONG:}}
		\begin{itemize}
			\item \textcolor{myred}{A:} PCA REDUCES features
			\item \textcolor{myred}{C:} PCA is not classification
			\item \textcolor{myred}{D:} PCA is not clustering
		\end{itemize}
		
		\textcolor{myred}{\textbf{EXAM TRAP:}} PCA = \textcolor{myblue}{dimensionality reduction, not classification/clustering}!
	\end{frame}
	
	\begin{frame}{Q73: Principal Components Definition - Tutorial}
		\fontsize{6.5}{7.5}\selectfont
		\textbf{QUESTION: What is the first principal component?}
		
		\vspace{0.1cm}
		\textbf{OPTIONS:}
		\begin{enumerate}
			\item The linear combination that captures the most variability in the data
			\item The average of all features
			\item The first feature in the dataset
			\item The feature with the highest value
		\end{enumerate}
		
		\vspace{0.1cm}
		\textcolor{myblue}{\textbf{ANSWER: A}}
		
		\vspace{0.1cm}
		\textbf{DETAILED EXPLANATION:}
		\begin{itemize}
			\item Linear combination of original predictors
			\item Captures most variability among all possible combinations
			\item Maximizes variance
			\item Component weights show feature importance
		\end{itemize}
		
		\vspace{0.1cm}
		\textcolor{myred}{\textbf{WHY OTHER OPTIONS ARE WRONG:}}
		\begin{itemize}
			\item \textcolor{myred}{B:} Average is not a principal component
			\item \textcolor{myred}{C:} First PC is a combination, not a single feature
			\item \textcolor{myred}{D:} PC is not about values, about variability
		\end{itemize}
		
		\textcolor{myred}{\textbf{EXAM TRAP:}} First PC = \textcolor{myblue}{captures most variability}!
	\end{frame}
	
	\begin{frame}{Q74: PC Orthogonality - Tutorial}
		\fontsize{6.5}{7.5}\selectfont
		\textbf{QUESTION: What does it mean that principal components are orthogonal?}
		
		\vspace{0.1cm}
		\textbf{OPTIONS:}
		\begin{enumerate}
			\item They are highly correlated
			\item They are uncorrelated (independent)
			\item They have equal variance
			\item They have equal means
		\end{enumerate}
		
		\vspace{0.1cm}
		\textcolor{myblue}{\textbf{ANSWER: B}}
		
		\vspace{0.1cm}
		\textbf{DETAILED EXPLANATION:}
		\begin{itemize}
			\item Orthogonal = uncorrelated
			\item Each PC captures different aspect of variability
			\item No overlap in information
			\item Useful for avoiding multicollinearity
		\end{itemize}
		
		\vspace{0.1cm}
		\textcolor{myred}{\textbf{WHY OTHER OPTIONS ARE WRONG:}}
		\begin{itemize}
			\item \textcolor{myred}{A:} Orthogonal means NOT correlated
			\item \textcolor{myred}{C:} PCs have decreasing variance
			\item \textcolor{myred}{D:} Means are not necessarily equal
		\end{itemize}
		
		\textcolor{myred}{\textbf{EXAM TRAP:}} Orthogonal = \textcolor{myblue}{uncorrelated}!
	\end{frame}
	
	\begin{frame}{Q75: PCA Formula - Tutorial}
		\fontsize{6.5}{7.5}\selectfont
		\textbf{QUESTION: What is the formula for a principal component?}
		
		\vspace{0.1cm}
		\textbf{OPTIONS:}
		\begin{enumerate}
			\item $PC_j = b_{1j}P_1 + b_{2j}P_2 + ... + b_{mj}P_m$
			\item $PC_j = \frac{1}{m}\sum_{i=1}^m P_i$
			\item $PC_j = \max(P_1, P_2, ..., P_m)$
			\item $PC_j = P_1 \times P_2 \times ... \times P_m$
		\end{enumerate}
		
		\vspace{0.1cm}
		\textcolor{myblue}{\textbf{ANSWER: A}}
		
		\vspace{0.1cm}
		\textbf{DETAILED EXPLANATION:}
		\begin{itemize}
			\item $P_i$: original predictors (standardized)
			\item $b_{ij}$: component weights
			\item Linear combination of original features
			\item Weights show importance of each feature
		\end{itemize}
		
		\vspace{0.1cm}
		\textcolor{myred}{\textbf{WHY OTHER OPTIONS ARE WRONG:}}
		\begin{itemize}
			\item \textcolor{myred}{B:} This is the mean
			\item \textcolor{myred}{C:} This is the maximum
			\item \textcolor{myred}{D:} This is the product
		\end{itemize}
		
		\textcolor{myred}{\textbf{EXAM TRAP:}} PC = \textcolor{myblue}{linear combination of features with weights}!
	\end{frame}
	
	\begin{frame}{Q76: Component Weights - Tutorial}
		\fontsize{6.5}{7.5}\selectfont
		\textbf{QUESTION: What do component weights ($b_{ij}$) represent in PCA?}
		
		\vspace{0.1cm}
		\textbf{OPTIONS:}
		\begin{enumerate}
			\item The mean of each feature
			\item The importance of each feature for the principal component
			\item The standard deviation of each feature
			\item The correlation between features
		\end{enumerate}
		
		\vspace{0.1cm}
		\textcolor{myblue}{\textbf{ANSWER: B}}
		
		\vspace{0.1cm}
		\textbf{DETAILED EXPLANATION:}
		\begin{itemize}
			\item Weights indicate feature importance
			\item Larger absolute weight = more important
			\item Positive weights: feature moves with PC
			\item Negative weights: feature moves opposite to PC
		\end{itemize}
		
		\vspace{0.1cm}
		\textcolor{myred}{\textbf{WHY OTHER OPTIONS ARE WRONG:}}
		\begin{itemize}
			\item \textcolor{myred}{A:} Means are different
			\item \textcolor{myred}{C:} Standard deviations are different
			\item \textcolor{myred}{D:} Correlations are different
		\end{itemize}
		
		\textcolor{myred}{\textbf{EXAM TRAP:}} Weights = \textcolor{myblue}{feature importance for each PC}!
	\end{frame}
	
	\begin{frame}{Q77: PCA Yield Curve Example - Tutorial}
		\fontsize{6.5}{7.5}\selectfont
		\textbf{QUESTION: In the Treasury yield curve PCA, what does PC1 represent?}
		
		\vspace{0.1cm}
		\textbf{OPTIONS:}
		\begin{enumerate}
			\item Slope of the yield curve
			\item Level of the yield curve
			\item Twist of the yield curve
			\item Curvature of the yield curve
		\end{enumerate}
		
		\vspace{0.1cm}
		\textcolor{myblue}{\textbf{ANSWER: B}}
		
		\vspace{0.1cm}
		\textbf{DETAILED EXPLANATION:}
		\begin{itemize}
			\item PC1 loads positively on all maturities
			\item All yields move in same direction
			\item Represents parallel shift (level)
			\item Explains 82\% of variation
		\end{itemize}
		
		\vspace{0.1cm}
		\textcolor{myred}{\textbf{WHY OTHER OPTIONS ARE WRONG:}}
		\begin{itemize}
			\item \textcolor{myred}{A:} Slope is PC2 (short/long opposite)
			\item \textcolor{myred}{C:} Twist is PC3
			\item \textcolor{myred}{D:} Curvature is not mentioned
		\end{itemize}
		
		\textcolor{myred}{\textbf{EXAM TRAP:}} PC1 = \textcolor{myblue}{level of yield curve}!
	\end{frame}
	
	\begin{frame}{Q78: PCA Yield Curve - PC2 - Tutorial}
		\fontsize{6.5}{7.5}\selectfont
		\textbf{QUESTION: In the Treasury yield curve PCA, what does PC2 represent?}
		
		\vspace{0.1cm}
		\textbf{OPTIONS:}
		\begin{enumerate}
			\item Level of the yield curve
			\item Slope of the yield curve
			\item Twist of the yield curve
			\item Curvature of the yield curve
		\end{enumerate}
		
		\vspace{0.1cm}
		\textcolor{myblue}{\textbf{ANSWER: B}}
		
		\vspace{0.1cm}
		\textbf{DETAILED EXPLANATION:}
		\begin{itemize}
			\item PC2 loads negatively on short-term yields
			\item Loads positively on long-term yields
			\item Represents steepening/flattening (slope)
			\item Short and long yields move opposite
		\end{itemize}
		
		\vspace{0.1cm}
		\textcolor{myred}{\textbf{WHY OTHER OPTIONS ARE WRONG:}}
		\begin{itemize}
			\item \textcolor{myred}{A:} Level is PC1
			\item \textcolor{myred}{C:} Twist is PC3
			\item \textcolor{myred}{D:} Curvature is not mentioned
		\end{itemize}
		
		\textcolor{myred}{\textbf{EXAM TRAP:}} PC2 = \textcolor{myblue}{slope of yield curve}!
	\end{frame}
	
	\begin{frame}{Q79: PCA Yield Curve - PC3 - Tutorial}
		\fontsize{6.5}{7.5}\selectfont
		\textbf{QUESTION: In the Treasury yield curve PCA, what does PC3 represent?}
		
		\vspace{0.1cm}
		\textbf{OPTIONS:}
		\begin{enumerate}
			\item Level of the yield curve
			\item Slope of the yield curve
			\item Twist of the yield curve
			\item The overall average
		\end{enumerate}
		
		\vspace{0.1cm}
		\textcolor{myblue}{\textbf{ANSWER: C}}
		
		\vspace{0.1cm}
		\textbf{DETAILED EXPLANATION:}
		\begin{itemize}
			\item PC3 loads positively on 1-year and 20+ year
			\item Loads negatively on 2-10 year yields
			\item Represents butterfly/twist
			\item Captures non-parallel shifts
		\end{itemize}
		
		\vspace{0.1cm}
		\textcolor{myred}{\textbf{WHY OTHER OPTIONS ARE WRONG:}}
		\begin{itemize}
			\item \textcolor{myred}{A:} Level is PC1
			\item \textcolor{myred}{B:} Slope is PC2
			\item \textcolor{myred}{D:} Not the average
		\end{itemize}
		
		\textcolor{myred}{\textbf{EXAM TRAP:}} PC3 = \textcolor{myblue}{twist of yield curve}!
	\end{frame}
	
	\begin{frame}{Q80: Variance Explained by PCs - Tutorial}
		\fontsize{6.5}{7.5}\selectfont
		\textbf{QUESTION: In the yield curve example, how much variance do the first 3 PCs explain?}
		
		\vspace{0.1cm}
		\textbf{OPTIONS:}
		\begin{enumerate}
			\item 50\%
			\item 75\%
			\item More than 98\%
			\item 100\%
		\end{enumerate}
		
		\vspace{0.1cm}
		\textcolor{myblue}{\textbf{ANSWER: C}}
		
		\vspace{0.1cm}
		\textbf{DETAILED EXPLANATION:}
		\begin{itemize}
			\item PC1: 82\%
			\item PC2: $\sim$12\% (to reach 94\%)
			\item PC3: $\sim$4\% (to reach 98+\%)
			\item Total variance explained by first 3 PCs $>$ 98\%
		\end{itemize}
		
		\vspace{0.1cm}
		\textcolor{myred}{\textbf{WHY OTHER OPTIONS ARE WRONG:}}
		\begin{itemize}
			\item \textcolor{myred}{A:} More than 50\%
			\item \textcolor{myred}{B:} More than 75\%
			\item \textcolor{myred}{D:} Not 100\% (some in later PCs)
		\end{itemize}
		
		\textcolor{myred}{\textbf{EXAM TRAP:}} First 3 PCs explain \textcolor{myblue}{$>$98\% of variance} in yield curve!
	\end{frame}
	
	\begin{frame}{Q81: PCA Elbow Plot - Tutorial}
		\fontsize{6.5}{7.5}\selectfont
		\textbf{QUESTION: How is the number of PCs chosen using an elbow plot?}
		
		\vspace{0.1cm}
		\textbf{OPTIONS:}
		\begin{enumerate}
			\item Choose all PCs
			\item Choose the point where explained variance drops off (elbow)
			\item Always choose the first 5 PCs
			\item Choose the PC with highest variance
		\end{enumerate}
		
		\vspace{0.1cm}
		\textcolor{myblue}{\textbf{ANSWER: B}}
		
		\vspace{0.1cm}
		\textbf{DETAILED EXPLANATION:}
		\begin{itemize}
			\item Plot: PC number vs variance explained
			\item Elbow: point where marginal variance drops
			\item Keep PCs before elbow
			\item After elbow: diminishing returns
		\end{itemize}
		
		\vspace{0.1cm}
		\textcolor{myred}{\textbf{WHY OTHER OPTIONS ARE WRONG:}}
		\begin{itemize}
			\item \textcolor{myred}{A:} Too many PCs may be unnecessary
			\item \textcolor{myred}{C:} Optimal number varies by dataset
			\item \textcolor{myred}{D:} Need to consider multiple PCs
		\end{itemize}
		
		\textcolor{myred}{\textbf{EXAM TRAP:}} Elbow = \textcolor{myblue}{point where variance drops off}!
	\end{frame}
	
	\begin{frame}{Q82: PCA and Multicollinearity - Tutorial}
		\fontsize{6.5}{7.5}\selectfont
		\textbf{QUESTION: How does PCA help with multicollinearity?}
		
		\vspace{0.1cm}
		\textbf{OPTIONS:}
		\begin{enumerate}
			\item It creates correlated features
			\item It creates uncorrelated features (orthogonal PCs)
			\item It removes all features
			\item It increases the number of features
		\end{enumerate}
		
		\vspace{0.1cm}
		\textcolor{myblue}{\textbf{ANSWER: B}}
		
		\vspace{0.1cm}
		\textbf{DETAILED EXPLANATION:}
		\begin{itemize}
			\item Original features may be highly correlated
			\item PCs are constructed to be uncorrelated
			\item Reduces multicollinearity problems
			\item Improves numerical stability of models
		\end{itemize}
		
		\vspace{0.1cm}
		\textcolor{myred}{\textbf{WHY OTHER OPTIONS ARE WRONG:}}
		\begin{itemize}
			\item \textcolor{myred}{A:} PCs are uncorrelated, not correlated
			\item \textcolor{myred}{C:} PCA doesn't remove all features
			\item \textcolor{myred}{D:} PCA reduces features
		\end{itemize}
		
		\textcolor{myred}{\textbf{EXAM TRAP:}} PCA creates \textcolor{myblue}{uncorrelated features} (solves multicollinearity)!
	\end{frame}
	
	\begin{frame}{Q83: PCA Requirements - Tutorial}
		\fontsize{6.5}{7.5}\selectfont
		\textbf{QUESTION: What should be done before applying PCA?}
		
		\vspace{0.1cm}
		\textbf{OPTIONS:}
		\begin{enumerate}
			\item Nothing - PCA handles everything
			\item Scale/standardize the features
			\item Remove all categorical variables
			\item Add more features
		\end{enumerate}
		
		\vspace{0.1cm}
		\textcolor{myblue}{\textbf{ANSWER: B}}
		
		\vspace{0.1cm}
		\textbf{DETAILED EXPLANATION:}
		\begin{itemize}
			\item Features should be scaled (mean 0, variance 1)
			\item Prevents features with large scales from dominating
			\item Remove skewness before scaling
			\item PCA is sensitive to scale differences
		\end{itemize}
		
		\vspace{0.1cm}
		\textcolor{myred}{\textbf{WHY OTHER OPTIONS ARE WRONG:}}
		\begin{itemize}
			\item \textcolor{myred}{A:} Scaling is necessary
			\item \textcolor{myred}{C:} Categorical variables need encoding first
			\item \textcolor{myred}{D:} Adding features defeats the purpose
		\end{itemize}
		
		\textcolor{myred}{\textbf{EXAM TRAP:}} \textcolor{myblue}{Scale/standardize} features before PCA!
	\end{frame}
	
	\begin{frame}{Q84: PCA Use Case - Bank of Jamaica - Tutorial}
		\fontsize{6.5}{7.5}\selectfont
		\textbf{QUESTION: How did Bank of Jamaica use PCA?}
		
		\vspace{0.1cm}
		\textbf{OPTIONS:}
		\begin{enumerate}
			\item For customer segmentation
			\item To develop a VaR methodology for interest rate risk
			\item For fraud detection
			\item For image recognition
		\end{enumerate}
		
		\vspace{0.1cm}
		\textcolor{myblue}{\textbf{ANSWER: B}}
		
		\vspace{0.1cm}
		\textbf{DETAILED EXPLANATION:}
		\begin{itemize}
			\item Developed PCA-based VaR methodology
			\item Assessed interest rate risk of Jamaican banks
			\item Decomposed yield curve movements into orthogonal factors
			\item More comprehensive than traditional VaR methods
		\end{itemize}
		
		\vspace{0.1cm}
		\textcolor{myred}{\textbf{WHY OTHER OPTIONS ARE WRONG:}}
		\begin{itemize}
			\item \textcolor{myred}{A:} This is a different application
			\item \textcolor{myred}{C:} Fraud detection is different
			\item \textcolor{myred}{D:} Image recognition is different
		\end{itemize}
		
		\textcolor{myred}{\textbf{EXAM TRAP:}} BOJ used PCA for \textcolor{myblue}{interest rate risk VaR}!
	\end{frame}
	
	\begin{frame}{Q85: PCA Summary - Tutorial}
		\fontsize{6.5}{7.5}\selectfont
		\textbf{QUESTION: What is the most comprehensive summary of PCA?}
		
		\vspace{0.1cm}
		\textbf{OPTIONS:}
		\begin{enumerate}
			\item PCA is a supervised learning method for classification
			\item PCA is an unsupervised technique that creates uncorrelated linear combinations of features, capturing variability, reducing dimensionality, and addressing multicollinearity
			\item PCA is a clustering algorithm
			\item PCA only works for categorical data
		\end{enumerate}
		
		\vspace{0.1cm}
		\textcolor{myblue}{\textbf{ANSWER: B}}
		
		\vspace{0.1cm}
		\textbf{DETAILED EXPLANATION:}
		\begin{itemize}
			\item \textbf{Unsupervised:} No labels needed
			\item \textbf{Linear combinations:} Creates new features
			\item \textbf{Uncorrelated:} PCs are orthogonal
			\item \textbf{Captures variability:} Ordered by variance explained
			\item \textbf{Reduces dimensionality:} Fewer features
			\item \textbf{Solves multicollinearity:} PCs are independent
		\end{itemize}
		
		\vspace{0.1cm}
		\textcolor{myred}{\textbf{WHY OTHER OPTIONS ARE WRONG:}}
		\begin{itemize}
			\item \textcolor{myred}{A:} PCA is UNSUPERVISED
			\item \textcolor{myred}{C:} PCA is for dimensionality reduction, not clustering
			\item \textcolor{myred}{D:} PCA is for numerical data
		\end{itemize}
		
		\textcolor{myred}{\textbf{EXAM TRAP:}} PCA = \textcolor{myblue}{unsupervised dimensionality reduction with uncorrelated components}!
	\end{frame}
	
	% ============================================================
	% SECTION 7: TRAINING, VALIDATION, TESTING - QUESTIONS 86-100
	% ============================================================
	
	\begin{frame}{Q86: Training Set Purpose - Tutorial}
		\fontsize{6.5}{7.5}\selectfont
		\textbf{QUESTION: What is the purpose of the training set?}
		
		\vspace{0.1cm}
		\textbf{OPTIONS:}
		\begin{enumerate}
			\item To select between competing models
			\item To evaluate final model performance
			\item To estimate model parameters (learning)
			\item To tune hyperparameters
		\end{enumerate}
		
		\vspace{0.1cm}
		\textcolor{myblue}{\textbf{ANSWER: C}}
		
		\vspace{0.1cm}
		\textbf{DETAILED EXPLANATION:}
		\begin{itemize}
			\item Used to estimate parameters (weights)
			\item The model "learns" from this data
			\item Largest portion of data typically
			\item Should represent the population
		\end{itemize}
		
		\vspace{0.1cm}
		\textcolor{myred}{\textbf{WHY OTHER OPTIONS ARE WRONG:}}
		\begin{itemize}
			\item \textcolor{myred}{A:} Validation set is for model selection
			\item \textcolor{myred}{B:} Test set is for final evaluation
			\item \textcolor{myred}{D:} Validation set is for hyperparameter tuning
		\end{itemize}
		
		\textcolor{myred}{\textbf{EXAM TRAP:}} Training set = \textcolor{myblue}{estimating parameters (learning)}!
	\end{frame}
	
	\begin{frame}{Q87: Validation Set Purpose - Tutorial}
		\fontsize{6.5}{7.5}\selectfont
		\textbf{QUESTION: What is the purpose of the validation set?}
		
		\vspace{0.1cm}
		\textbf{OPTIONS:}
		\begin{enumerate}
			\item To estimate model parameters
			\item To select between competing models and tune hyperparameters
			\item To test final model performance
			\item To increase data size
		\end{enumerate}
		
		\vspace{0.1cm}
		\textcolor{myblue}{\textbf{ANSWER: B}}
		
		\vspace{0.1cm}
		\textbf{DETAILED EXPLANATION:}
		\begin{itemize}
			\item Used for model selection
			\item Tune hyperparameters
			\item Compare competing models
			\item Once used, it's "contaminated"
		\end{itemize}
		
		\vspace{0.1cm}
		\textcolor{myred}{\textbf{WHY OTHER OPTIONS ARE WRONG:}}
		\begin{itemize}
			\item \textcolor{myred}{A:} Training set estimates parameters
			\item \textcolor{myred}{C:} Test set evaluates final performance
			\item \textcolor{myred}{D:} Validation set doesn't increase data
		\end{itemize}
		
		\textcolor{myred}{\textbf{EXAM TRAP:}} Validation set = \textcolor{myblue}{model selection and hyperparameter tuning}!
	\end{frame}
	
	\begin{frame}{Q88: Test Set Purpose - Tutorial}
		\fontsize{6.5}{7.5}\selectfont
		\textbf{QUESTION: What is the purpose of the test set?}
		
		\vspace{0.1cm}
		\textbf{OPTIONS:}
		\begin{enumerate}
			\item To estimate model parameters
			\item To select between competing models
			\item To evaluate the final model's performance on unseen data
			\item To tune hyperparameters
		\end{enumerate}
		
		\vspace{0.1cm}
		\textcolor{myblue}{\textbf{ANSWER: C}}
		
		\vspace{0.1cm}
		\textbf{DETAILED EXPLANATION:}
		\begin{itemize}
			\item Held back until final evaluation
			\item Assess generalization performance
			\item Truly unseen data
			\item Provides unbiased performance estimate
		\end{itemize}
		
		\vspace{0.1cm}
		\textcolor{myred}{\textbf{WHY OTHER OPTIONS ARE WRONG:}}
		\begin{itemize}
			\item \textcolor{myred}{A:} Training set estimates parameters
			\item \textcolor{myred}{B:} Validation set selects models
			\item \textcolor{myred}{D:} Validation set tunes hyperparameters
		\end{itemize}
		
		\textcolor{myred}{\textbf{EXAM TRAP:}} Test set = \textcolor{myblue}{final unbiased evaluation}!
	\end{frame}
	
	\begin{frame}{Q89: Overfitting Definition - Tutorial}
		\fontsize{6.5}{7.5}\selectfont
		\textbf{QUESTION: What is overfitting in ML?}
		
		\vspace{0.1cm}
		\textbf{OPTIONS:}
		\begin{enumerate}
			\item Model fits training data well but fails to generalize
			\item Model fits test data better than training data
			\item Model is too simple to capture patterns
			\item Model has no parameters
		\end{enumerate}
		
		\vspace{0.1cm}
		\textcolor{myblue}{\textbf{ANSWER: A}}
		
		\vspace{0.1cm}
		\textbf{DETAILED EXPLANATION:}
		\begin{itemize}
			\item Model captures noise, not signal
			\item Excellent training performance
			\item Poor test performance
			\item Too complex model
		\end{itemize}
		
		\vspace{0.1cm}
		\textcolor{myred}{\textbf{WHY OTHER OPTIONS ARE WRONG:}}
		\begin{itemize}
			\item \textcolor{myred}{B:} This would be unusual and suspicious
			\item \textcolor{myred}{C:} This is underfitting
			\item \textcolor{myred}{D:} Overfitting involves too many parameters
		\end{itemize}
		
		\textcolor{myred}{\textbf{EXAM TRAP:}} Overfitting = \textcolor{myblue}{good training, poor test}!
	\end{frame}
	
	\begin{frame}{Q90: Underfitting Definition - Tutorial}
		\fontsize{6.5}{7.5}\selectfont
		\textbf{QUESTION: What is underfitting in ML?}
		
		\vspace{0.1cm}
		\textbf{OPTIONS:}
		\begin{enumerate}
			\item Model fits training data well but fails to generalize
			\item Model is too simple to capture patterns
			\item Model fits test data perfectly
			\item Model has too many parameters
		\end{enumerate}
		
		\vspace{0.1cm}
		\textcolor{myblue}{\textbf{ANSWER: B}}
		
		\vspace{0.1cm}
		\textbf{DETAILED EXPLANATION:}
		\begin{itemize}
			\item Too simple model
			\item Cannot capture relationships
			\item Poor performance on both training and test
			\item Opposite of overfitting
		\end{itemize}
		
		\vspace{0.1cm}
		\textcolor{myred}{\textbf{WHY OTHER OPTIONS ARE WRONG:}}
		\begin{itemize}
			\item \textcolor{myred}{A:} This is overfitting
			\item \textcolor{myred}{C:} This would be suspicious
			\item \textcolor{myred}{D:} This describes overfitting
		\end{itemize}
		
		\textcolor{myred}{\textbf{EXAM TRAP:}} Underfitting = \textcolor{myblue}{model too simple}!
	\end{frame}
	
	\begin{frame}{Q91: Sample Split Ratio - Tutorial}
		\fontsize{6.5}{7.5}\selectfont
		\textbf{QUESTION: What is a common rule of thumb for sample splitting?}
		
		\vspace{0.1cm}
		\textbf{OPTIONS:}
		\begin{enumerate}
			\item Training: 50\%, Validation: 25\%, Test: 25\%
			\item Training: 66\%, Validation: 17\%, Test: 17\%
			\item Training: 80\%, Validation: 10\%, Test: 10\%
			\item Training: 90\%, Validation: 5\%, Test: 5\%
		\end{enumerate}
		
		\vspace{0.1cm}
		\textcolor{myblue}{\textbf{ANSWER: B}}
		
		\vspace{0.1cm}
		\textbf{DETAILED EXPLANATION:}
		\begin{itemize}
			\item Roughly two-thirds training
			\item Remaining third split equally between validation and test
			\item Approximate guideline
			\item May vary depending on data size
		\end{itemize}
		
		\vspace{0.1cm}
		\textcolor{myred}{\textbf{WHY OTHER OPTIONS ARE WRONG:}}
		\begin{itemize}
			\item \textcolor{myred}{A:} Training is too small
			\item \textcolor{myred}{C:} Validation/test too small
			\item \textcolor{myred}{D:} Validation/test too small
		\end{itemize}
		
		\textcolor{myred}{\textbf{EXAM TRAP:}} Common split = \textcolor{myblue}{66\% training, 17\% validation, 17\% test}!
	\end{frame}
	
	\begin{frame}{Q92: Cross-Sectional Sample Split - Tutorial}
		\fontsize{6.5}{7.5}\selectfont
		\textbf{QUESTION: How should cross-sectional data be split?}
		
		\vspace{0.1cm}
		\textbf{OPTIONS:}
		\begin{enumerate}
			\item Chronologically (first part training, last part test)
			\item Randomly from the total dataset
			\item Always use all data for training
			\item Sort by feature values
		\end{enumerate}
		
		\vspace{0.1cm}
		\textcolor{myblue}{\textbf{ANSWER: B}}
		
		\vspace{0.1cm}
		\textbf{DETAILED EXPLANATION:}
		\begin{itemize}
			\item Cross-sectional data has no order
			\item Random sampling ensures representativeness
			\item Avoids bias
			\item Each observation independent
		\end{itemize}
		
		\vspace{0.1cm}
		\textcolor{myred}{\textbf{WHY OTHER OPTIONS ARE WRONG:}}
		\begin{itemize}
			\item \textcolor{myred}{A:} This is for time-series data
			\item \textcolor{myred}{C:} Need to hold out data for testing
			\item \textcolor{myred}{D:} Should be random, not sorted
		\end{itemize}
		
		\textcolor{myred}{\textbf{EXAM TRAP:}} Cross-sectional = \textcolor{myblue}{random split}!
	\end{frame}
	
	\begin{frame}{Q93: Time-Series Sample Split - Tutorial}
		\fontsize{6.5}{7.5}\selectfont
		\textbf{QUESTION: How should time-series data be split?}
		
		\vspace{0.1cm}
		\textbf{OPTIONS:}
		\begin{enumerate}
			\item Randomly from the total dataset
			\item Chronologically: first part training, last part test
			\item Always use all data for training
			\item Sort by feature values
		\end{enumerate}
		
		\vspace{0.1cm}
		\textcolor{myblue}{\textbf{ANSWER: B}}
		
		\vspace{0.1cm}
		\textbf{DETAILED EXPLANATION:}
		\begin{itemize}
			\item Time-series has chronological order
			\item Training = earlier periods
			\item Validation = middle periods
			\item Test = most recent periods
			\item Avoids look-ahead bias
		\end{itemize}
		
		\vspace{0.1cm}
		\textcolor{myred}{\textbf{WHY OTHER OPTIONS ARE WRONG:}}
		\begin{itemize}
			\item \textcolor{myred}{A:} Random split would break chronology
			\item \textcolor{myred}{C:} Need to hold out data for testing
			\item \textcolor{myred}{D:} Chronology matters, not feature values
		\end{itemize}
		
		\textcolor{myred}{\textbf{EXAM TRAP:}} Time-series = \textcolor{myblue}{chronological split}!
	\end{frame}
	
	\begin{frame}{Q94: Cross-Validation Introduction - Tutorial}
		\fontsize{6.5}{7.5}\selectfont
		\textbf{QUESTION: When is cross-validation used?}
		
		\vspace{0.1cm}
		\textbf{OPTIONS:}
		\begin{enumerate}
			\item When there is abundant data
			\item When there is limited data
			\item Always
			\item Never
		\end{enumerate}
		
		\vspace{0.1cm}
		\textcolor{myblue}{\textbf{ANSWER: B}}
		
		\vspace{0.1cm}
		\textbf{DETAILED EXPLANATION:}
		\begin{itemize}
			\item Used when data is limited
			\item Combines training and validation
			\item Test data still held back
			\item More efficient use of data
		\end{itemize}
		
		\vspace{0.1cm}
		\textcolor{myred}{\textbf{WHY OTHER OPTIONS ARE WRONG:}}
		\begin{itemize}
			\item \textcolor{myred}{A:} With abundant data, simple split works
			\item \textcolor{myred}{C:} Only needed for limited data
			\item \textcolor{myred}{D:} Cross-validation is often used
		\end{itemize}
		
		\textcolor{myred}{\textbf{EXAM TRAP:}} Cross-validation = \textcolor{myblue}{for limited data}!
	\end{frame}
	
	\begin{frame}{Q95: K-Fold Cross-Validation - Tutorial}
		\fontsize{6.5}{7.5}\selectfont
		\textbf{QUESTION: How does k-fold cross-validation work?}
		
		\vspace{0.1cm}
		\textbf{OPTIONS:}
		\begin{enumerate}
			\item Use all data for training once
			\item Split data into k folds, train on k-1, validate on 1, repeat k times
			\item Split data once into training and test
			\item Use all data for validation
		\end{enumerate}
		
		\vspace{0.1cm}
		\textcolor{myblue}{\textbf{ANSWER: B}}
		
		\vspace{0.1cm}
		\textbf{DETAILED EXPLANATION:}
		\begin{itemize}
			\item Combines training and validation
			\item k folds (e.g., 5 or 10)
			\item Each fold used as validation once
			\item Results averaged across folds
		\end{itemize}
		
		\vspace{0.1cm}
		\textcolor{myred}{\textbf{WHY OTHER OPTIONS ARE WRONG:}}
		\begin{itemize}
			\item \textcolor{myred}{A:} This is just training
			\item \textcolor{myred}{C:} This is simple split
			\item \textcolor{myred}{D:} This is no training
		\end{itemize}
		
		\textcolor{myred}{\textbf{EXAM TRAP:}} k-fold CV = \textcolor{myblue}{train on k-1, validate on 1, repeat k times}!
	\end{frame}
	
	\begin{frame}{Q96: Good Generalization - Tutorial}
		\fontsize{6.5}{7.5}\selectfont
		\textbf{QUESTION: What indicates a model generalizes well?}
		
		\vspace{0.1cm}
		\textbf{OPTIONS:}
		\begin{enumerate}
			\item Training accuracy is much higher than test accuracy
			\item Test accuracy is similar to training accuracy
			\item Training accuracy is very low
			\item Test accuracy is very low
		\end{enumerate}
		
		\vspace{0.1cm}
		\textcolor{myblue}{\textbf{ANSWER: B}}
		
		\vspace{0.1cm}
		\textbf{DETAILED EXPLANATION:}
		\begin{itemize}
			\item Similar performance on training and test
			\item Model captures signal, not noise
			\item Not overfitting or underfitting
			\item Good predictive performance
		\end{itemize}
		
		\vspace{0.1cm}
		\textcolor{myred}{\textbf{WHY OTHER OPTIONS ARE WRONG:}}
		\begin{itemize}
			\item \textcolor{myred}{A:} This indicates overfitting
			\item \textcolor{myred}{C:} This indicates underfitting
			\item \textcolor{myred}{D:} This indicates poor performance
		\end{itemize}
		
		\textcolor{myred}{\textbf{EXAM TRAP:}} Good generalization = \textcolor{myblue}{similar training and test performance}!
	\end{frame}
	
	\begin{frame}{Q97: Model Selection Process - Tutorial}
		\fontsize{6.5}{7.5}\selectfont
		\textbf{QUESTION: What is the correct order of model development?}
		
		\vspace{0.1cm}
		\textbf{OPTIONS:}
		\begin{enumerate}
			\item Train $\to$ Test $\to$ Validate
			\item Train $\to$ Validate $\to$ Test
			\item Validate $\to$ Train $\to$ Test
			\item Test $\to$ Train $\to$ Validate
		\end{enumerate}
		
		\vspace{0.1cm}
		\textcolor{myblue}{\textbf{ANSWER: B}}
		
		\vspace{0.1cm}
		\textbf{DETAILED EXPLANATION:}
		\begin{enumerate}
			\item Train: Estimate parameters
			\item Validate: Select model and tune hyperparameters
			\item Test: Evaluate final model performance
			\item Test data used only once at the end
		\end{enumerate}
		
		\vspace{0.1cm}
		\textcolor{myred}{\textbf{WHY OTHER OPTIONS ARE WRONG:}}
		\begin{itemize}
			\item \textcolor{myred}{A:} Test should be last
			\item \textcolor{myred}{C:} Training must come before validation
			\item \textcolor{myred}{D:} All steps are out of order
		\end{itemize}
		
		\textcolor{myred}{\textbf{EXAM TRAP:}} Correct order = \textcolor{myblue}{Train $\to$ Validate $\to$ Test}!
	\end{frame}
	
	\begin{frame}{Q98: Software for ML - Tutorial}
		\fontsize{6.5}{7.5}\selectfont
		\textbf{QUESTION: Which programming languages are popular for ML?}
		
		\vspace{0.1cm}
		\textbf{OPTIONS:}
		\begin{enumerate}
			\item R and Python
			\item Java and C++
			\item HTML and CSS
			\item SQL and PHP
		\end{enumerate}
		
		\vspace{0.1cm}
		\textcolor{myblue}{\textbf{ANSWER: A}}
		
		\vspace{0.1cm}
		\textbf{DETAILED EXPLANATION:}
		\begin{itemize}
			\item \textbf{R:} Great for statistics and visualization
			\item \textbf{Python:} General-purpose, great for ML
			\item Both open-source, cross-platform
			\item Rich ecosystem of libraries
		\end{itemize}
		
		\vspace{0.1cm}
		\textcolor{myred}{\textbf{WHY OTHER OPTIONS ARE WRONG:}}
		\begin{itemize}
			\item \textcolor{myred}{B:} Java/C++ are less common for ML
			\item \textcolor{myred}{C:} HTML/CSS are web languages
			\item \textcolor{myred}{D:} SQL/PHP are database/web languages
		\end{itemize}
		
		\textcolor{myred}{\textbf{EXAM TRAP:}} Popular ML languages = \textcolor{myblue}{R and Python}!
	\end{frame}
	
	\begin{frame}{Q99: Python Libraries - Tutorial}
		\fontsize{6.5}{7.5}\selectfont
		\textbf{QUESTION: Which Python library is for machine learning?}
		
		\vspace{0.1cm}
		\textbf{OPTIONS:}
		\begin{enumerate}
			\item NumPy
			\item Scikit-Learn
			\item Matplotlib
			\item All of the above
		\end{enumerate}
		
		\vspace{0.1cm}
		\textcolor{myblue}{\textbf{ANSWER: D}}
		
		\vspace{0.1cm}
		\textbf{DETAILED EXPLANATION:}
		\begin{itemize}
			\item \textbf{NumPy:} Array operations, linear algebra
			\item \textbf{Scikit-Learn:} ML models, preprocessing
			\item \textbf{Matplotlib:} Data visualization
			\item All part of Python ML ecosystem
		\end{itemize}
		
		\vspace{0.1cm}
		\textcolor{myred}{\textbf{WHY OTHER OPTIONS ARE WRONG:}}
		\begin{itemize}
			\item \textcolor{myred}{A:} Correct but incomplete
			\item \textcolor{myred}{B:} Correct but incomplete
			\item \textcolor{myred}{C:} Correct but incomplete
		\end{itemize}
		
		\textcolor{myred}{\textbf{EXAM TRAP:}} All are \textcolor{myblue}{important Python libraries} for ML!
	\end{frame}
	
	\begin{frame}{Q100: Final Summary - Tutorial}
		\fontsize{6.5}{7.5}\selectfont
		\textbf{QUESTION: What is the most comprehensive summary of ML fundamentals?}
		
		\vspace{0.1cm}
		\textbf{OPTIONS:}
		\begin{enumerate}
			\item ML is only about supervised learning
			\item ML includes supervised, unsupervised, semi-supervised, and reinforcement learning, with data pre-processing, PCA, and proper train/validation/test splits
			\item ML only uses Python
			\item ML doesn't require data cleaning
		\end{enumerate}
		
		\vspace{0.1cm}
		\textcolor{myblue}{\textbf{ANSWER: B}}
		
		\vspace{0.1cm}
		\textbf{DETAILED EXPLANATION:}
		\begin{itemize}
			\item \textbf{ML Types:} Supervised, Unsupervised, Semi-supervised, Reinforcement
			\item \textbf{Data Pre-processing:} Cleaning, scaling, encoding
			\item \textbf{PCA:} Dimensionality reduction
			\item \textbf{Splitting:} Training, Validation, Test
		\end{itemize}
		
		\vspace{0.1cm}
		\textcolor{myred}{\textbf{WHY OTHER OPTIONS ARE WRONG:}}
		\begin{itemize}
			\item \textcolor{myred}{A:} ML includes more than supervised
			\item \textcolor{myred}{C:} R is also used
			\item \textcolor{myred}{D:} Data cleaning is essential
		\end{itemize}
		
		\textcolor{myred}{\textbf{EXAM SUCCESS:}} Understand \textcolor{myblue}{ML types, data pre-processing, PCA, and data splitting}!
	\end{frame}
	
\end{document}